% =====================================================================================================================
% Revised Section 4 (Method). Draft for the author to edit; the original sections/4_method.tex is unchanged.
%
% How to use: in main.tex replace \section{Model and Estimation}
\label{sec:method}

\subsection{Base-model selection as a discrete choice}
Following \citet{overgoor2019choosing}, each first adoption $k$ is a choice of one parent $j_k$ from the market
$\M_k$. Developer $k$ values candidate $j$ at $U_{kj} = \theta^\top x_{kj} + \varepsilon_{kj}$ with i.i.d.\ Gumbel
errors, which gives the conditional logit \citep{mcfadden1974conditional}
\begin{equation}
  P(j \mid \M_k) = \frac{\exp(\theta^\top x_{kj})}{\sum_{l \in \M_k} \exp(\theta^\top x_{kl})}.
  \label{eq:logit}
\end{equation}
The features $x_{kj}$ are the four groups of \cref{sec:data}; they vary over candidates and over time and, for
demand match and inertia, over developers. With $\theta = 0$ the model is \emph{uniform attachment}, every
alive model equally likely; preferential attachment is the special case in which only $\log d$ and $\1[d>0]$
enter. We interpret $\theta$ descriptively, as the weights that best summarize which signals choices co-move with,
not as causal effects (\cref{sec:discussion}).

\subsection{Sampling the competitors that matter}
\label{sec:sampling}
Evaluating \eqref{eq:logit} over a median of \marketMedian{} alternatives for \nFirst{} choices is infeasible, so we
sample a candidate set $\D_k \subset \M_k$. Each alternative enters independently (Poisson sampling) with
\begin{equation}
  \pi_{kl} = \min\{1,\; c_k (d_{kl} + 1)\}, \qquad \textstyle\sum_{l \in \M_k} \pi_{kl} = S,
  \label{eq:pps}
\end{equation}
where $d_{kl}$ is the degree of $l$ just before choice $k$ and $c_k$ is solved in closed form from the degree
histogram of the market; the chosen parent is always included. We use $S=300$; on average 8.6 candidates per choice
have $\pi=1$, so the incumbents a developer is most likely to weigh are present in every set, and the rest of the
set is a probability-proportional-to-size sample of the long tail.

\begin{proposition}[Sampling correction]
\label{prop:pps}
Under \eqref{eq:logit} and the sampling scheme \eqref{eq:pps},
\begin{equation}
  P(j \mid \D_k) = \frac{\exp(\theta^\top x_{kj} - \log \pi_{kj})}{\sum_{l \in \D_k} \exp(\theta^\top x_{kl} - \log
  \pi_{kl})},
  \label{eq:corrected}
\end{equation}
and maximizing $\sum_k \log P(j_k \mid \D_k)$ gives a consistent estimator of $\theta$.
\end{proposition}
The proof (\cref{app:sampling}) applies McFadden's positive-conditioning argument \citep{mcfadden1978modelling}:
given that $j$ was chosen, the probability of drawing $\D$ is $\prod_{l \in \D \setminus j} \pi_l
\prod_{l \notin \D}(1-\pi_l)$, so the ratio for two members of $\D$ is $\pi_l / \pi_j$. The offset $-\log \pi$
makes uniform attachment, $\theta = 0$, the natural baseline: it assigns every sampled candidate the weight
$1/\pi$, i.e.\ every model in the full market the same probability.

Why sample proportionally to degree? Consistency holds for any $\pi > 0$, but efficiency does not. A uniform sample
of 300 from a market of $3 \times 10^5$ almost never contains the dozen models that receive most adoptions, so each
choice compares the chosen parent mostly with never-used models, and the likelihood learns little about how
developers trade off one popular model against another. \Cref{sec:robustq} shows that this is not only a matter of
precision: with uniform sets the estimated weights change.

\subsection{Estimation and inference}
The training data have 35.2 million (choice, candidate) rows. We group choices into buckets of similar set size,
keep the 26 features on the GPU in single precision and accumulate the log-likelihood, gradient and Hessian in double
precision; L-BFGS with an analytic Hessian fits the full model in under a minute on one A100. Standard errors are
clustered by author \citep{cameron2015practitioner}: \nFirstAuthors{} authors make the \nFirst{} first adoptions, and
the choices of one author are strongly dependent (\cref{app:validation}). For every derived quantity, including the
Shapley shares below, we use a multiplier (Bayesian) bootstrap over authors \citep{rubin1981bayesian}: each
replicate draws a weight $w_a \sim \mathrm{Exp}(1)$ per author, reweights the author's choices and refits from the
point estimate.

\subsection{Decomposing the likelihood gain}
\label{sec:shapley}
The four groups are correlated (popular models also tend to be well documented and to match common tasks), so the
change in fit from adding a group depends on which groups are already in the model. We therefore attribute the
log-likelihood gain over uniform attachment, $G(\mathcal{S}) = \ell(\hat\theta_{\mathcal{S}}) - \ell(0)$ for the
model with the groups in $\mathcal{S}$, by the Shapley value \citep{shapley1953value}
\begin{equation}
  \phi_g = \sum_{\mathcal{S} \subseteq \mathcal{G} \setminus \{g\}}
  \frac{|\mathcal{S}|!\,(|\mathcal{G}|-|\mathcal{S}|-1)!}{|\mathcal{G}|!}
  \big[G(\mathcal{S} \cup \{g\}) - G(\mathcal{S})\big],
  \label{eq:shapley}
\end{equation}
and report shares $\phi_g / G(\mathcal{G})$. With four groups this needs $2^4 = 16$ fits, each warm-started from
a subset one group smaller. We compute the shares both in-sample and on the holdout (mean per-choice
log-likelihood), where a group can only gain by predicting future choices.

\subsection{Bounding general quality}
\label{sec:bounds}
Benchmark scores cover few models, so a small quality share could reflect missing measurement rather than
indifference to quality. We bound the quality share from above by adding an intercept $\phi_j$ for each of the
500 most adopted parents (assigned ex ante, \cref{sec:data}). The fixed effects absorb every time-invariant
attribute of a hub: unmeasured quality, but also brand, tutorials and default status. Entering the fixed effects as
a fifth Shapley player and counting their entire share as quality gives the largest share that any fixed hub
attribute could contribute. To see how much of this bound is quality, we regress $\hat\phi_j$ on external quality
measures (Open LLM Leaderboard v2 average, arena ratings; controls for parameter count and re-distributor accounts)
and correct the $R^2$ for estimation noise using the bootstrap variance of $\hat\phi_j$; the noise-corrected share
of $\mathrm{Var}(\phi)$ explained by quality splits the fixed-effect share into a quality and a popularity part.

\subsection{Does the decomposition recover known mechanisms?}
\label{sec:validation}
Two checks support the procedure (details in \cref{app:validation}). First, on synthetic preferential-attachment
networks the corrected likelihood recovers the attachment exponent under all three sampling designs. Second, we
replace every observed choice with a draw from a known utility on the same candidate sets and re-run the full
decomposition. When choices are drawn from the popularity terms alone, the decomposition attributes
\recSfPop\% of the gain to popularity; when drawn from quality and demand match alone, it attributes \recQfFit\% to
fit; when drawn from the full estimated model, it returns \recFullPop\%, \recFullMatch\% and \recFullInertia\% for
popularity, demand match and inertia, the shares estimated on the real choices. The decomposition can tell the
mechanisms apart, and it does not manufacture a popularity share where none was generated.
 by % =====================================================================================================================
% Revised Section 4 (Method). Draft for the author to edit; the original sections/4_method.tex is unchanged.
%
% How to use: in main.tex replace \section{Model and Estimation}
\label{sec:method}

\subsection{Base-model selection as a discrete choice}
Following \citet{overgoor2019choosing}, each first adoption $k$ is a choice of one parent $j_k$ from the market
$\M_k$. Developer $k$ values candidate $j$ at $U_{kj} = \theta^\top x_{kj} + \varepsilon_{kj}$ with i.i.d.\ Gumbel
errors, which gives the conditional logit \citep{mcfadden1974conditional}
\begin{equation}
  P(j \mid \M_k) = \frac{\exp(\theta^\top x_{kj})}{\sum_{l \in \M_k} \exp(\theta^\top x_{kl})}.
  \label{eq:logit}
\end{equation}
The features $x_{kj}$ are the four groups of \cref{sec:data}; they vary over candidates and over time and, for
demand match and inertia, over developers. With $\theta = 0$ the model is \emph{uniform attachment}, every
alive model equally likely; preferential attachment is the special case in which only $\log d$ and $\1[d>0]$
enter. We interpret $\theta$ descriptively, as the weights that best summarize which signals choices co-move with,
not as causal effects (\cref{sec:discussion}).

\subsection{Sampling the competitors that matter}
\label{sec:sampling}
Evaluating \eqref{eq:logit} over a median of \marketMedian{} alternatives for \nFirst{} choices is infeasible, so we
sample a candidate set $\D_k \subset \M_k$. Each alternative enters independently (Poisson sampling) with
\begin{equation}
  \pi_{kl} = \min\{1,\; c_k (d_{kl} + 1)\}, \qquad \textstyle\sum_{l \in \M_k} \pi_{kl} = S,
  \label{eq:pps}
\end{equation}
where $d_{kl}$ is the degree of $l$ just before choice $k$ and $c_k$ is solved in closed form from the degree
histogram of the market; the chosen parent is always included. We use $S=300$; on average 8.6 candidates per choice
have $\pi=1$, so the incumbents a developer is most likely to weigh are present in every set, and the rest of the
set is a probability-proportional-to-size sample of the long tail.

\begin{proposition}[Sampling correction]
\label{prop:pps}
Under \eqref{eq:logit} and the sampling scheme \eqref{eq:pps},
\begin{equation}
  P(j \mid \D_k) = \frac{\exp(\theta^\top x_{kj} - \log \pi_{kj})}{\sum_{l \in \D_k} \exp(\theta^\top x_{kl} - \log
  \pi_{kl})},
  \label{eq:corrected}
\end{equation}
and maximizing $\sum_k \log P(j_k \mid \D_k)$ gives a consistent estimator of $\theta$.
\end{proposition}
The proof (\cref{app:sampling}) applies McFadden's positive-conditioning argument \citep{mcfadden1978modelling}:
given that $j$ was chosen, the probability of drawing $\D$ is $\prod_{l \in \D \setminus j} \pi_l
\prod_{l \notin \D}(1-\pi_l)$, so the ratio for two members of $\D$ is $\pi_l / \pi_j$. The offset $-\log \pi$
makes uniform attachment, $\theta = 0$, the natural baseline: it assigns every sampled candidate the weight
$1/\pi$, i.e.\ every model in the full market the same probability.

Why sample proportionally to degree? Consistency holds for any $\pi > 0$, but efficiency does not. A uniform sample
of 300 from a market of $3 \times 10^5$ almost never contains the dozen models that receive most adoptions, so each
choice compares the chosen parent mostly with never-used models, and the likelihood learns little about how
developers trade off one popular model against another. \Cref{sec:robustq} shows that this is not only a matter of
precision: with uniform sets the estimated weights change.

\subsection{Estimation and inference}
The training data have 35.2 million (choice, candidate) rows. We group choices into buckets of similar set size,
keep the 26 features on the GPU in single precision and accumulate the log-likelihood, gradient and Hessian in double
precision; L-BFGS with an analytic Hessian fits the full model in under a minute on one A100. Standard errors are
clustered by author \citep{cameron2015practitioner}: \nFirstAuthors{} authors make the \nFirst{} first adoptions, and
the choices of one author are strongly dependent (\cref{app:validation}). For every derived quantity, including the
Shapley shares below, we use a multiplier (Bayesian) bootstrap over authors \citep{rubin1981bayesian}: each
replicate draws a weight $w_a \sim \mathrm{Exp}(1)$ per author, reweights the author's choices and refits from the
point estimate.

\subsection{Decomposing the likelihood gain}
\label{sec:shapley}
The four groups are correlated (popular models also tend to be well documented and to match common tasks), so the
change in fit from adding a group depends on which groups are already in the model. We therefore attribute the
log-likelihood gain over uniform attachment, $G(\mathcal{S}) = \ell(\hat\theta_{\mathcal{S}}) - \ell(0)$ for the
model with the groups in $\mathcal{S}$, by the Shapley value \citep{shapley1953value}
\begin{equation}
  \phi_g = \sum_{\mathcal{S} \subseteq \mathcal{G} \setminus \{g\}}
  \frac{|\mathcal{S}|!\,(|\mathcal{G}|-|\mathcal{S}|-1)!}{|\mathcal{G}|!}
  \big[G(\mathcal{S} \cup \{g\}) - G(\mathcal{S})\big],
  \label{eq:shapley}
\end{equation}
and report shares $\phi_g / G(\mathcal{G})$. With four groups this needs $2^4 = 16$ fits, each warm-started from
a subset one group smaller. We compute the shares both in-sample and on the holdout (mean per-choice
log-likelihood), where a group can only gain by predicting future choices.

\subsection{Bounding general quality}
\label{sec:bounds}
Benchmark scores cover few models, so a small quality share could reflect missing measurement rather than
indifference to quality. We bound the quality share from above by adding an intercept $\phi_j$ for each of the
500 most adopted parents (assigned ex ante, \cref{sec:data}). The fixed effects absorb every time-invariant
attribute of a hub: unmeasured quality, but also brand, tutorials and default status. Entering the fixed effects as
a fifth Shapley player and counting their entire share as quality gives the largest share that any fixed hub
attribute could contribute. To see how much of this bound is quality, we regress $\hat\phi_j$ on external quality
measures (Open LLM Leaderboard v2 average, arena ratings; controls for parameter count and re-distributor accounts)
and correct the $R^2$ for estimation noise using the bootstrap variance of $\hat\phi_j$; the noise-corrected share
of $\mathrm{Var}(\phi)$ explained by quality splits the fixed-effect share into a quality and a popularity part.

\subsection{Does the decomposition recover known mechanisms?}
\label{sec:validation}
Two checks support the procedure (details in \cref{app:validation}). First, on synthetic preferential-attachment
networks the corrected likelihood recovers the attachment exponent under all three sampling designs. Second, we
replace every observed choice with a draw from a known utility on the same candidate sets and re-run the full
decomposition. When choices are drawn from the popularity terms alone, the decomposition attributes
\recSfPop\% of the gain to popularity; when drawn from quality and demand match alone, it attributes \recQfFit\% to
fit; when drawn from the full estimated model, it returns \recFullPop\%, \recFullMatch\% and \recFullInertia\% for
popularity, demand match and inertia, the shares estimated on the real choices. The decomposition can tell the
mechanisms apart, and it does not manufacture a popularity share where none was generated.
 by % =====================================================================================================================
% Revised Section 4 (Method). Draft for the author to edit; the original sections/4_method.tex is unchanged.
%
% How to use: in main.tex replace \section{Model and Estimation}
\label{sec:method}

\subsection{Base-model selection as a discrete choice}
Following \citet{overgoor2019choosing}, each first adoption $k$ is a choice of one parent $j_k$ from the market
$\M_k$. Developer $k$ values candidate $j$ at $U_{kj} = \theta^\top x_{kj} + \varepsilon_{kj}$ with i.i.d.\ Gumbel
errors, which gives the conditional logit \citep{mcfadden1974conditional}
\begin{equation}
  P(j \mid \M_k) = \frac{\exp(\theta^\top x_{kj})}{\sum_{l \in \M_k} \exp(\theta^\top x_{kl})}.
  \label{eq:logit}
\end{equation}
The features $x_{kj}$ are the four groups of \cref{sec:data}; they vary over candidates and over time and, for
demand match and inertia, over developers. With $\theta = 0$ the model is \emph{uniform attachment}, every
alive model equally likely; preferential attachment is the special case in which only $\log d$ and $\1[d>0]$
enter. We interpret $\theta$ descriptively, as the weights that best summarize which signals choices co-move with,
not as causal effects (\cref{sec:discussion}).

\subsection{Sampling the competitors that matter}
\label{sec:sampling}
Evaluating \eqref{eq:logit} over a median of \marketMedian{} alternatives for \nFirst{} choices is infeasible, so we
sample a candidate set $\D_k \subset \M_k$. Each alternative enters independently (Poisson sampling) with
\begin{equation}
  \pi_{kl} = \min\{1,\; c_k (d_{kl} + 1)\}, \qquad \textstyle\sum_{l \in \M_k} \pi_{kl} = S,
  \label{eq:pps}
\end{equation}
where $d_{kl}$ is the degree of $l$ just before choice $k$ and $c_k$ is solved in closed form from the degree
histogram of the market; the chosen parent is always included. We use $S=300$; on average 8.6 candidates per choice
have $\pi=1$, so the incumbents a developer is most likely to weigh are present in every set, and the rest of the
set is a probability-proportional-to-size sample of the long tail.

\begin{proposition}[Sampling correction]
\label{prop:pps}
Under \eqref{eq:logit} and the sampling scheme \eqref{eq:pps},
\begin{equation}
  P(j \mid \D_k) = \frac{\exp(\theta^\top x_{kj} - \log \pi_{kj})}{\sum_{l \in \D_k} \exp(\theta^\top x_{kl} - \log
  \pi_{kl})},
  \label{eq:corrected}
\end{equation}
and maximizing $\sum_k \log P(j_k \mid \D_k)$ gives a consistent estimator of $\theta$.
\end{proposition}
The proof (\cref{app:sampling}) applies McFadden's positive-conditioning argument \citep{mcfadden1978modelling}:
given that $j$ was chosen, the probability of drawing $\D$ is $\prod_{l \in \D \setminus j} \pi_l
\prod_{l \notin \D}(1-\pi_l)$, so the ratio for two members of $\D$ is $\pi_l / \pi_j$. The offset $-\log \pi$
makes uniform attachment, $\theta = 0$, the natural baseline: it assigns every sampled candidate the weight
$1/\pi$, i.e.\ every model in the full market the same probability.

Why sample proportionally to degree? Consistency holds for any $\pi > 0$, but efficiency does not. A uniform sample
of 300 from a market of $3 \times 10^5$ almost never contains the dozen models that receive most adoptions, so each
choice compares the chosen parent mostly with never-used models, and the likelihood learns little about how
developers trade off one popular model against another. \Cref{sec:robustq} shows that this is not only a matter of
precision: with uniform sets the estimated weights change.

\subsection{Estimation and inference}
The training data have 35.2 million (choice, candidate) rows. We group choices into buckets of similar set size,
keep the 26 features on the GPU in single precision and accumulate the log-likelihood, gradient and Hessian in double
precision; L-BFGS with an analytic Hessian fits the full model in under a minute on one A100. Standard errors are
clustered by author \citep{cameron2015practitioner}: \nFirstAuthors{} authors make the \nFirst{} first adoptions, and
the choices of one author are strongly dependent (\cref{app:validation}). For every derived quantity, including the
Shapley shares below, we use a multiplier (Bayesian) bootstrap over authors \citep{rubin1981bayesian}: each
replicate draws a weight $w_a \sim \mathrm{Exp}(1)$ per author, reweights the author's choices and refits from the
point estimate.

\subsection{Decomposing the likelihood gain}
\label{sec:shapley}
The four groups are correlated (popular models also tend to be well documented and to match common tasks), so the
change in fit from adding a group depends on which groups are already in the model. We therefore attribute the
log-likelihood gain over uniform attachment, $G(\mathcal{S}) = \ell(\hat\theta_{\mathcal{S}}) - \ell(0)$ for the
model with the groups in $\mathcal{S}$, by the Shapley value \citep{shapley1953value}
\begin{equation}
  \phi_g = \sum_{\mathcal{S} \subseteq \mathcal{G} \setminus \{g\}}
  \frac{|\mathcal{S}|!\,(|\mathcal{G}|-|\mathcal{S}|-1)!}{|\mathcal{G}|!}
  \big[G(\mathcal{S} \cup \{g\}) - G(\mathcal{S})\big],
  \label{eq:shapley}
\end{equation}
and report shares $\phi_g / G(\mathcal{G})$. With four groups this needs $2^4 = 16$ fits, each warm-started from
a subset one group smaller. We compute the shares both in-sample and on the holdout (mean per-choice
log-likelihood), where a group can only gain by predicting future choices.

\subsection{Bounding general quality}
\label{sec:bounds}
Benchmark scores cover few models, so a small quality share could reflect missing measurement rather than
indifference to quality. We bound the quality share from above by adding an intercept $\phi_j$ for each of the
500 most adopted parents (assigned ex ante, \cref{sec:data}). The fixed effects absorb every time-invariant
attribute of a hub: unmeasured quality, but also brand, tutorials and default status. Entering the fixed effects as
a fifth Shapley player and counting their entire share as quality gives the largest share that any fixed hub
attribute could contribute. To see how much of this bound is quality, we regress $\hat\phi_j$ on external quality
measures (Open LLM Leaderboard v2 average, arena ratings; controls for parameter count and re-distributor accounts)
and correct the $R^2$ for estimation noise using the bootstrap variance of $\hat\phi_j$; the noise-corrected share
of $\mathrm{Var}(\phi)$ explained by quality splits the fixed-effect share into a quality and a popularity part.

\subsection{Does the decomposition recover known mechanisms?}
\label{sec:validation}
Two checks support the procedure (details in \cref{app:validation}). First, on synthetic preferential-attachment
networks the corrected likelihood recovers the attachment exponent under all three sampling designs. Second, we
replace every observed choice with a draw from a known utility on the same candidate sets and re-run the full
decomposition. When choices are drawn from the popularity terms alone, the decomposition attributes
\recSfPop\% of the gain to popularity; when drawn from quality and demand match alone, it attributes \recQfFit\% to
fit; when drawn from the full estimated model, it returns \recFullPop\%, \recFullMatch\% and \recFullInertia\% for
popularity, demand match and inertia, the shares estimated on the real choices. The decomposition can tell the
mechanisms apart, and it does not manufacture a popularity share where none was generated.
 by % =====================================================================================================================
% Revised Section 4 (Method). Draft for the author to edit; the original sections/4_method.tex is unchanged.
%
% How to use: in main.tex replace \input{sections/4_method} by \input{sections/4_method_revised}, and add
% \input{appendix/method_revised_proofs} next to the other appendices (it holds the proofs cited here).
% All labels used elsewhere in the paper (sec:method, eq:logit, sec:sampling, eq:pps, prop:pps, eq:corrected,
% sec:shapley, eq:shapley, sec:bounds, sec:validation) are kept.
%
% Theorem environments: acmart defines definition, lemma, proposition, corollary; assumption and remark are
% defined below if missing (no change to main.tex needed).
%
% Changes relative to 4_method.tex, by comment of the author:
%  1. 4.1 now defines the choice occasion, degree and the information set (Def. 1) before the random-utility
%     assumption (Assumption 1); IIA is stated, because Prop. 1 uses it.
%  2. Consistency / asymptotic normality is a separate corollary (Cor. 1) after the sampling proposition.
%  3. 4.3 is about the estimator: weighted sampled likelihood, gradient and Hessian, concavity (Lemma 3), the
%     estimation algorithm (Algorithm 1), the clustered sandwich and the multiplier bootstrap. Hardware is one
%     sentence.
%  4. 4.4 starts from why coefficients cannot answer the question (units, groups, within-set variation, order
%     dependence with the real numbers) and defines the likelihood-gain game (Def. 3) before the Shapley value.
%  5. 4.5 defines the hub set formally (Def. 4) and states the noise-corrected split as a lemma.
%  6. 4.6 writes the three generating utilities, the simulated choice and the recovery criterion as equations.
%  New 4.7: samples, weights and author types (definitions only; results in Section 5, Table tab:samples).
%  Five groups (self-supply and habit replace inertia) and 2^5 = 32 fits throughout.
%
% Section 3 needs two matching edits (not done): (a) "the 500 most adopted parents" -> refer to Def. 4 (the hubs are
% ranked by degree d_l(t_h) over all children, not by first adoptions); (b) the inertia group -> self-supply
% (same_author) and habit (author_reuse, author_family, is_last).
% =====================================================================================================================
\makeatletter
\@ifundefined{assumption}{\newtheorem{assumption}{Assumption}}{}
\@ifundefined{remark}{\theoremstyle{acmdefinition}\newtheorem{remark}[theorem]{Remark}\theoremstyle{acmplain}}{}
\makeatother

\section{Model and Estimation}
\label{sec:method}

\Cref{sec:data} showed that adoption is highly concentrated and that popularity, general quality and demand match
could each produce this concentration. To separate them we need a model in which all three can act at once and be
compared on a common scale. This section builds that model (\cref{sec:choice}), makes it computable on markets of
hundreds of thousands of alternatives (\cref{sec:sampling,sec:estimation}), turns its estimates into a
decomposition by signal group (\cref{sec:shapley}), bounds what unmeasured quality could add (\cref{sec:bounds}),
and checks that the decomposition recovers known mechanisms (\cref{sec:validation}).

% ---------------------------------------------------------------------------------------------------------------------
\subsection{Base-model selection as a discrete choice}
\label{sec:choice}

\begin{definition}[Choice occasion]
\label{def:occasion}
A \emph{choice occasion} $k$ is a first adoption (\cref{sec:data}), described by its author $a_k$, time $t_k$,
modality $m_k$ and chosen parent $j_k$. Its \emph{market} $\M_k$ is the set of models of modality $m_k$ created
before $t_k$ that $a_k$ has not used as a parent. The \emph{degree} of a model $l$ at time $t$ is
\[
  d_l(t) = \#\{c : \text{$c$ declares $l$ as its parent with type fine-tune, adapter or unspecified, and $c$ was
  created before $t$}\},
\]
and $d_{kl} = d_l(t_k)$. Every candidate $l \in \M_k$ carries a feature vector $x_{kl} \in \mathbb{R}^p$ computed
only from models, edges and metadata that exist before $t_k$.
\end{definition}

The features are the five groups $g \in \mathcal G = \{\text{pop}, \text{qual}, \text{match}, \text{self},
\text{habit}\}$ of \cref{sec:data}: popularity, general quality, demand match, self-supply (the candidate is the
author's own model) and habit (the candidate belongs to a family the author used before). We write $x_{kl} =
(x_{kl,g})_{g \in \mathcal G}$ and $\theta = (\theta_g)_{g \in \mathcal G}$ conformably. Popularity and quality
features vary over candidates and time; demand-match, self-supply and habit features also vary over authors.

\begin{assumption}[Random utility]
\label{ass:rum}
The author of occasion $k$ chooses $j_k = \arg\max_{l \in \M_k} U_{kl}$, where
\[
  U_{kl} = V_{kl} + \varepsilon_{kl}, \qquad V_{kl} = \theta^\top x_{kl},
\]
and the $\varepsilon_{kl}$ are i.i.d.\ standard Type-I extreme value, independent of $\{x_{kl}\}_{l \in \M_k}$.
\end{assumption}

\begin{lemma}[\citealp{mcfadden1974conditional}]
\label{lem:logit}
Under \cref{ass:rum},
\begin{equation}
  P(j \mid \M_k) = \frac{\exp(\theta^\top x_{kj})}{\sum_{l \in \M_k} \exp(\theta^\top x_{kl})}.
  \label{eq:logit}
\end{equation}
\end{lemma}

\begin{remark}[What the model does and does not say]
\label{rem:logit}
(i) \emph{Identification.} Only differences across candidates within an occasion enter \eqref{eq:logit}; a
variable that is constant across $l \in \M_k$ (author experience, organizational status) cancels. Author
heterogeneity therefore enters through author-specific features (demand match, self-supply, habit) and through
estimation by author type (\cref{sec:samples}).
(ii) \emph{Scale.} The error variance is normalized, so $\theta$ is measured relative to the dispersion of tastes;
comparisons across samples are made with likelihood shares (\cref{sec:shapley}), which are scale free, not with
coefficients.
(iii) \emph{Independence of irrelevant alternatives.} $P(j \mid \M_k) / P(l \mid \M_k) = \exp(V_{kj} - V_{kl})$
does not depend on the other alternatives. \Cref{prop:pps} uses this property.
(iv) \emph{Interpretation.} We read $\theta$ descriptively, as the weights that best summarize which signals
choices co-move with, not as causal effects (\cref{sec:discussion}).
\end{remark}

Two network-growth models are special cases. With $\theta = 0$, \eqref{eq:logit} is \emph{uniform attachment}:
every model in the market is equally likely. With only $x_{kl} = (\log \max(d_{kl},1),\ \1[d_{kl}>0])$ and
coefficients $(\alpha, \beta)$,
\[
  P(j \mid \M_k) \;\propto\;
  \begin{cases} e^{\beta}\, d_{kj}^{\alpha}, & d_{kj} > 0, \\ 1, & d_{kj} = 0, \end{cases}
\]
which is \emph{preferential attachment} with exponent $\alpha$; $\alpha = 1$ is the linear model of
\citet{barabasi1999emergence}. Following \citet{overgoor2019choosing}, \eqref{eq:logit} nests these mechanisms and
lets them be compared by likelihood together with quality and demand match.

% ---------------------------------------------------------------------------------------------------------------------
\subsection{Sampling the competitors that matter}
\label{sec:sampling}

The denominator of \eqref{eq:logit} has a median of \marketMedian{} terms, and the training data contain \nTrain{}
occasions, so the likelihood cannot be evaluated on full markets. We evaluate it on a sampled candidate set.

\begin{definition}[Degree-proportional Poisson design]
\label{def:pps}
For occasion $k$ and a target size $S$, every $l \in \M_k$ enters the candidate set $\D_k$ independently with
probability
\begin{equation}
  \pi_{kl} = \min\{1,\; c_k (d_{kl} + 1)\},
  \label{eq:pps}
\end{equation}
where the scale $c_k > 0$ solves $\sum_{l \in \M_k} \pi_{kl} = S$. The chosen parent $j_k$ is added to $\D_k$ if it
was not drawn.
\end{definition}

The inclusion probability is proportional to $d+1$ (the $+1$ keeps never-used models in the design), and $c_k$ is
the factor of proportionality. The condition on $c_k$ fixes the expected size of the set:

\begin{lemma}[Set size and the scale $c_k$]
\label{lem:size}
(i) $\E|\D_k| = \sum_{l \in \M_k} \pi_{kl} + (1 - \pi_{kj_k}) \in [S, S+1]$.
(ii) Let $n_\delta$ be the number of models in $\M_k$ with degree $\delta$, and let $\delta_1 > \delta_2 > \dots >
\delta_M > 0$ be the positive degrees. The function $\Sigma(c) = \sum_\delta n_\delta \min\{1, c(\delta+1)\}$ is
continuous and non-decreasing, strictly increasing wherever some degree class is unsaturated
($c(\delta+1) < 1$), and runs from $0$ to $|\M_k|$. If $|\M_k| > S$, $\Sigma(c) = S$ has a unique solution
\[
  c_k = \frac{S - \sum_{i \le m} n_{\delta_i}}{n_0 + \sum_{i > m} n_{\delta_i}(\delta_i + 1)},
\]
where $m \in \{0, \dots, M\}$ is the number of saturated classes: the unique $m$ with $c_k(\delta_i + 1) \ge 1$ for
$i \le m$, $c_k(\delta_i + 1) < 1$ for $i > m$, and $c_k \le 1$.
\end{lemma}

\noindent The proof is in \cref{app:proofs}. With $S = 300$, on average \piOneMean{} candidates per training
occasion have $\pi = 1$; these are the incumbents a developer is most likely to weigh, and \piOneChosen\% of chosen
parents are among them. The rest of each set is a probability-proportional-to-size sample of the long tail.

\begin{proposition}[Sampling correction]
\label{prop:pps}
Under \cref{ass:rum} and \cref{def:pps},
\begin{equation}
  P(j \mid \D_k) = \frac{\exp(\theta^\top x_{kj} - \log \pi_{kj})}{\sum_{l \in \D_k} \exp(\theta^\top x_{kl} - \log
  \pi_{kl})}, \qquad j \in \D_k.
  \label{eq:corrected}
\end{equation}
\end{proposition}

\noindent Given that $j$ was chosen, the probability of drawing a set $\D \ni j$ is $q(\D \mid j) = \prod_{l \in \D
\setminus \{j\}} \pi_l \prod_{l \notin \D} (1 - \pi_l)$, so $q(\D \mid j)/q(\D \mid l) = \pi_l / \pi_j$ for any $j,
l \in \D$; Bayes' rule and \eqref{eq:logit} give \eqref{eq:corrected} (\cref{app:proofs};
\citealp{mcfadden1978modelling}). The argument uses only the logit form of $P(j \mid \M_k)$, so it holds for any
specification of $V_{kl}$, including the fixed-effect model of \cref{sec:bounds}.

\begin{corollary}[Consistency and asymptotic normality]
\label{cor:consistency}
Let $\hat\theta$ maximize $\ell(\theta) = \sum_k \log P(j_k \mid \D_k; \theta)$ from \eqref{eq:corrected}. Suppose
occasions are independent across authors with bounded numbers per author, the parameter space is compact with the
true $\theta_0$ in its interior, $\E\|x_{kl}\|^2 < \infty$, and $H_0 = \E[\mathrm{Cov}_{p_k}(x_k)]$ is positive
definite, where $\mathrm{Cov}_{p_k}$ is the covariance of $x_{kl}$ over $l \in \D_k$ under the probabilities
\eqref{eq:corrected}. Then, as the number of authors $A \to \infty$, $\hat\theta \to \theta_0$ in probability and
$\sqrt{A}(\hat\theta - \theta_0) \to N(0, H_0^{-1} \Omega_0 H_0^{-1})$, with $\Omega_0$ the variance of an author's
summed score.
\end{corollary}

\noindent By \cref{prop:pps}, $\ell$ is a correctly specified conditional likelihood of the observed $(j_k, \D_k)$,
and it is concave (\cref{lem:concave} below); the result follows from the standard argument for concave M-estimators
\citep{newey1994large}. Two consequences of \eqref{eq:corrected} matter for what follows.

\begin{remark}[Uniform attachment as the baseline, and why sample by degree]
\label{rem:pps}
(i) At $\theta = 0$, \eqref{eq:corrected} gives each sampled candidate the weight $1/\pi_{kl}$, and $\sum_{l \in
\D_k} 1/\pi_{kl}$ is a Horvitz--Thompson estimate of $|\M_k|$ \citep{horvitz1952generalization}; uniform attachment
on the sample therefore corresponds to uniform attachment on the full market. It is the baseline of every likelihood
gain below.
(ii) \Cref{cor:consistency} holds for any design with $\pi > 0$, but the information matrix $H_0 = \E[\mathrm{Cov}_{p_k}
(x_k)]$ depends on the design: a candidate set contributes information only through the variation of $x$ among the
candidates that have non-negligible choice probability. A uniform sample of 300 from $3 \times 10^5$ models almost
never contains the few models that receive most adoptions, so the chosen parent is compared mostly with never-used
models. Under a uniform design the median clustered standard error is \seRatioUnif{} times larger, and the
estimated weights themselves change (\cref{sec:robustq}).
\end{remark}

% ---------------------------------------------------------------------------------------------------------------------
\subsection{Estimation and inference}
\label{sec:estimation}

\Cref{cor:consistency} defines the estimator; we now state how it is computed and how its uncertainty is measured.
Write $v_{kl}(\theta) = \theta^\top x_{kl} - \log \pi_{kl}$, $p_{kl}(\theta) = e^{v_{kl}} / \sum_{m \in \D_k}
e^{v_{km}}$ and $\bar x_k(\theta) = \sum_{l \in \D_k} p_{kl} x_{kl}$. With occasion weights $w_k$ ($w_k = 1$ in the
main specification; other weights in \cref{sec:samples} and the bootstrap below), the sampled log-likelihood is
\begin{equation}
  \ell(\theta) = \sum_k w_k \Big[ v_{kj_k}(\theta) - \log \sum_{l \in \D_k} e^{v_{kl}(\theta)} \Big].
  \label{eq:loglik}
\end{equation}

\begin{lemma}[Concavity]
\label{lem:concave}
$\nabla \ell(\theta) = \sum_k w_k (x_{kj_k} - \bar x_k)$ and $\nabla^2 \ell(\theta) = -\sum_k w_k
\mathrm{Cov}_{p_k}(x_k)$. Hence $\ell$ is concave, and strictly concave, with a unique maximizer, whenever
$\sum_k w_k \mathrm{Cov}_{p_k}(x_k)$ is positive definite.
\end{lemma}

\noindent The offset $-\log \pi$ does not depend on $\theta$, so it changes neither the gradient's form nor the
concavity. The first-order condition equates the weighted mean features of the chosen parents with their mean under
the model. Concavity has a practical consequence: every refit in the decomposition and the bootstrap converges to
the same optimum from any starting point, so warm starts change only computing time.

\begin{center}
\fbox{\begin{minipage}{0.94\linewidth}
\textbf{Algorithm 1 (estimation).}
\begin{enumerate}
\item \emph{Replay.} Process the network in time order; at each occasion $k$ update $d_l(t_k)$, solve $c_k$
  (\cref{lem:size}), draw $\D_k$ by \cref{def:pps}, and record $(x_{kl}, \log \pi_{kl})_{l \in \D_k}$.
\item \emph{Fit.} Maximize \eqref{eq:loglik} by L-BFGS with the analytic gradient of \cref{lem:concave}; compute
  the Hessian at $\hat\theta$.
\item \emph{Subsets.} For every $\mathcal S \subseteq \mathcal G$, maximize \eqref{eq:loglik} with $\theta_g = 0$ for
  $g \notin \mathcal S$, warm-started from the fit of $\mathcal S$ minus one group ($2^{|\mathcal G|} = \nShapFits{}$
  fits; \cref{sec:shapley}).
\item \emph{Bootstrap.} Repeat steps 2--3 $B = 50$ times with weights $w_k \xi_{a_k}$, $\xi_a \sim \mathrm{Exp}(1)$
  i.i.d.\ over authors, warm-started from the point estimates.
\end{enumerate}
\end{minipage}}
\end{center}

\noindent The training data have \nRowsTrain{} million (occasion, candidate) rows; the full model fits in about a
minute on one GPU.\footnote{Candidate sets are grouped into buckets of similar size and kept on the GPU in single
precision; the likelihood, gradient and Hessian are accumulated in double precision.}

\paragraph{Inference.}
Occasions of one author are dependent (the same project, the same tastes, batch uploads): the \nFirst{} first
adoptions come from \nFirstAuthors{} authors, and the median clustered standard error is \seClusterInfl{} times
the one that treats occasions as independent. We therefore use the cluster-robust sandwich
\[
  \widehat{\mathrm{Var}}(\hat\theta) = \hat H^{-1} \Big( \textstyle\sum_a \hat s_a \hat s_a^\top \Big) \hat H^{-1},
  \qquad \hat H = -\nabla^2 \ell(\hat\theta), \quad \hat s_a = \sum_{k: a_k = a} w_k (x_{kj_k} - \bar x_k(\hat\theta)),
\]
\citep{cameron2015practitioner}. For derived quantities, the Shapley shares of \cref{sec:shapley} in particular, we
use the multiplier bootstrap of step 4: replicate $b$ maximizes $\sum_k \xi^{(b)}_{a_k} w_k \ell_k(\theta)$; the
shares are recomputed in each replicate; and intervals are the 2.5\% and 97.5\% percentiles. Normalized Exp(1)
weights are Dirichlet(1,\dots,1), the Bayesian bootstrap of \citet{rubin1981bayesian}; reweighting whole authors
preserves within-author dependence, and the weighted bootstrap is valid for M-estimators of this kind
\citep{ma2005robust}.

\paragraph{Out-of-sample evaluation.}
Let $K_h$ be the holdout occasions (the last 10\% in time) and $P_0$ uniform attachment ($\theta = 0$ in
\eqref{eq:corrected}). We report the mean holdout log-likelihood gain
\[
  \Delta_h = \frac{1}{|K_h|} \sum_{k \in K_h} \big[ \log P(j_k \mid \D_k; \hat\theta) - \log P_0(j_k \mid \D_k)
  \big],
\]
so that $e^{\Delta_h}$ is the geometric-mean factor by which the model raises the probability of the true parent
over uniform attachment ($\Delta_h = \holdGain$, a factor of \holdGainTimes), and the top-$k$ rate, the share of
holdout occasions in which $j_k$ ranks among the $k$ candidates with the highest $\hat V_{kl}$ (ranking by
utility, without the offset; ties broken at random).

% ---------------------------------------------------------------------------------------------------------------------
\subsection{From coefficients to a decomposition}
\label{sec:shapley}

$\hat\theta$ has 26 components, but the question is about groups of signals: how much of developers' choices do
popularity signals account for, and how much do quality and demand match? The coefficients cannot answer it
directly, for four reasons.
(i) Features have different units (indicators, logarithms of counts, similarities), so coefficients are not
comparable across features.
(ii) Each group contains several features, and the question concerns the group.
(iii) What a group contributes to explaining a choice depends on how much its \emph{index}
$I_{g,kl} = \theta_g^\top x_{kl,g}$ varies across the candidates of $\D_k$, not on the size of $\theta_g$ alone;
a large coefficient on a feature that rarely varies within a set explains little.
(iv) The groups are correlated: popular models are also well documented and match common tasks. The fit gained by
adding a group therefore depends on which groups are already in the model. Demand match accounts for \matchFirst\%
of the full model's likelihood gain when it enters first and for \matchLast\% when it enters last.
We therefore measure each group by its contribution to the model's fit and allocate that contribution with a rule
that does not depend on the order of entry.

\begin{definition}[Likelihood-gain game]
\label{def:game}
For $\mathcal S \subseteq \mathcal G$, let $\hat\theta_{\mathcal S}$ maximize \eqref{eq:loglik} subject to
$\theta_g = 0$ for $g \notin \mathcal S$, and let
\[
  G(\mathcal S) = \ell(\hat\theta_{\mathcal S}) - \ell(0),
\]
the log-likelihood gain over uniform attachment of the model with the groups in $\mathcal S$. The pair
$(\mathcal G, G)$ is a cooperative game with $G(\emptyset) = 0$. Its holdout version $G_h(\mathcal S)$ evaluates the
same $\hat\theta_{\mathcal S}$ by the mean holdout log-likelihood gain.
\end{definition}

Because the models are nested and fitted by maximum likelihood, $G$ is monotone: $\mathcal S \subseteq \mathcal T$
implies $G(\mathcal S) \le G(\mathcal T)$. We allocate $G(\mathcal G)$ by the Shapley value
\citep{shapley1953value},
\begin{equation}
  \phi_g = \sum_{\mathcal S \subseteq \mathcal G \setminus \{g\}}
  \frac{|\mathcal S|!\,(|\mathcal G|-|\mathcal S|-1)!}{|\mathcal G|!}
  \big[G(\mathcal S \cup \{g\}) - G(\mathcal S)\big],
  \qquad s_g = \frac{\phi_g}{G(\mathcal G)},
  \label{eq:shapley}
\end{equation}
the average marginal gain of group $g$ over all orders of entry. It is the only allocation that is efficient
($\sum_g \phi_g = G(\mathcal G)$), symmetric, assigns zero to a group that never adds fit, and is additive in $G$; by
monotonicity, every $s_g \in [0,1]$. Shapley allocations of fit are standard for correlated regressors
\citep{lipovetsky2001analysis,huettner2012axiomatic}. To report coarser groups, such as \emph{fit} = quality $\cup$
demand match, we apply \eqref{eq:shapley} to the game on blocks $v(\mathcal T) = G(\cup \mathcal T)$, whose values
are among the \nShapFits{} fits already computed. Shares are reported in-sample and on the holdout; on the holdout a
group gains only by predicting future choices, and $G_h$ need not be monotone.

% ---------------------------------------------------------------------------------------------------------------------
\subsection{Bounding general quality}
\label{sec:bounds}

Benchmark scores cover few models, so a small quality share could reflect missing measurement rather than
indifference to quality. We bound it from above with model-specific intercepts for the most used models.

\begin{definition}[Hubs]
\label{def:hubs}
Let $t_h$ be the train--holdout cut (\tHold). The \emph{hub set} $\mathcal H$ is the 500 models with the highest
degree $d_l(t_h)$ over the whole Hub (all modalities; ties broken by model index). $\mathcal H$ uses only training
data and does not depend on the sampling design.
\end{definition}

The fixed-effect model adds an intercept for each hub,
\begin{equation}
  V^{\mathrm{FE}}_{kl} = \theta^\top x_{kl} + \sum_{h \in \mathcal H} \phi_h \1[l = h],
  \label{eq:fe}
\end{equation}
estimated by \eqref{eq:corrected} with the same offset (\cref{prop:pps} holds for any $V$) and a small ridge penalty
$\frac{\lambda}{2}\|\phi\|^2$, $\lambda = 10^{-4}$. Non-hubs have $\phi = 0$, so $\phi_h$ is hub $h$'s
time-invariant attractiveness relative to the rest of the market. It absorbs every fixed attribute of the hub:
unmeasured quality, but also brand, tutorials and default status. Entering the fixed effects as an additional player
in \cref{def:game} gives a share $s_{\mathrm{FE}}$, and
\[
  \bar s_{\mathrm{qual}} = s_{\mathrm{qual}} + s_{\mathrm{FE}}
\]
is the largest share that any time-invariant hub attribute could contribute to quality (\feShare\% of the gain
in-sample).

To see how much of the fixed effects is quality, we project the estimated intercepts on external quality measures
$z_h$ (Open LLM Leaderboard v2 average; arena ratings) with controls $w_h$ (parameter count, re-distributor accounts),
$\hat\phi_h = b_0 + b^\top z_h + \gamma^\top w_h + u_h$, over the \extHubs{} hubs that are covered and precisely
estimated.

\begin{lemma}[Noise-corrected quality share]
\label{lem:noise}
Suppose $\hat\phi_h = \phi_h + e_h$ with $\E[e_h \mid \phi_h, z_h, w_h] = 0$ and $\mathrm{Var}(e_h) = \sigma_h^2$.
Then $\mathrm{Var}(\hat\phi) = \mathrm{Var}(\phi) + \bar\sigma^2$, the projection coefficients are unaffected by $e$,
and
\[
  q = \frac{\mathrm{Var}(\hat b^\top z_h)}{\mathrm{Var}(\hat\phi_h) - \bar\sigma^2}
\]
estimates the share of $\mathrm{Var}(\phi)$ explained by external quality, with $\sigma_h^2$ the bootstrap variance
of $\hat\phi_h$.
\end{lemma}

\noindent We split the fixed-effect share proportionally, $s^{\mathrm{ext}}_{\mathrm{qual}} = s_{\mathrm{qual}} + q\,
s_{\mathrm{FE}}$ ($q = \extQ\%$). Shapley shares are not linear in a split of a player, so this split is an
approximation; the two end points, $s_{\mathrm{qual}}$ and $\bar s_{\mathrm{qual}}$, do not depend on it
(\cref{app:proofs}).

% ---------------------------------------------------------------------------------------------------------------------
\subsection{Does the decomposition recover known mechanisms?}
\label{sec:validation}

The decomposition attributes most of the likelihood gain to popularity (\cref{sec:results}). Because the groups are
correlated, this could in principle be a property of the method rather than of the data. We test the method on
choices whose generating mechanism is known by construction.

For a generating utility $V^{(r)}$, keep every real candidate set $\D_k$ with its real features and inclusion
probabilities, and replace the observed parent $j_k$ by a draw
\begin{equation}
  \tilde j^{(r)}_k \sim P^{(r)}(l \mid \D_k) = \frac{\exp(V^{(r)}_{kl} - \log \pi_{kl})}{\sum_{m \in \D_k}
  \exp(V^{(r)}_{km} - \log \pi_{km})}, \qquad l \in \D_k,
  \label{eq:sir}
\end{equation}
the model-implied choice on the sampled set (\cref{prop:pps}). We use three utilities built from the estimates of
\cref{sec:estimation,sec:bounds}:
\[
  V^{\mathrm{SF}}_{kl} = \hat\theta_{\mathrm{pop}}^\top x_{kl,\mathrm{pop}}, \qquad
  V^{\mathrm{QF}}_{kl} = \hat\theta_{\mathrm{qual}}^\top x_{kl,\mathrm{qual}} + \hat\theta_{\mathrm{match}}^\top
  x_{kl,\mathrm{match}} + \hat q_l, \qquad
  V^{\mathrm{FULL}}_{kl} = \hat V^{\mathrm{FE}}_{kl},
\]
where $\hat q_l = \hat b^\top z_l$ is the part of the hub intercept explained by external quality ($\hat q_l = 0$ for
non-hubs). SF lets choices depend only on popularity, QF only on quality and demand match, and FULL on the estimated
model. Running steps 2--3 of Algorithm 1 on $\{(\tilde j^{(r)}_k, \D_k)\}_k$ gives shares $\hat s(r)$. If the
decomposition separates the mechanisms, $\hat s_{\mathrm{pop}}(\mathrm{SF})$ and $\hat s_{\mathrm{fit}}(\mathrm{QF})$
should be close to one, and $\hat s(\mathrm{FULL})$ should equal the shares estimated on the real choices; $1 -
\hat s_{\mathrm{pop}}(\mathrm{SF})$ and $1 - \hat s_{\mathrm{fit}}(\mathrm{QF})$ measure how much fit leaks to a
correlated group.

The decomposition attributes \recSfPop\% of the gain to popularity under SF and \recQfFit\% to fit under QF, and it
returns \recFullPop\%, \recFullMatch\% and \recFullInertia\% (popularity, demand match, self-supply plus habit) under
FULL, the shares estimated on the real choices.\todo{Recovery is run with four groups; rerun with five or keep the
note that self-supply + habit = the former inertia group.} When choices are generated by quality and demand match
alone, the decomposition does not manufacture a popularity share. As a check of the whole pipeline, we also grow
synthetic networks in which a new node attaches to $l$ with weight $d_l^{\alpha}$ if $d_l > 0$ and $e^{\beta_0}$
otherwise ($\alpha = 1$, $\beta_0 = -1$), replay them with the same code, and recover $\hat\alpha$ within three
standard errors of one under all three sampling designs (\cref{app:validation}).

% ---------------------------------------------------------------------------------------------------------------------
\subsection{Samples, weights and author types}
\label{sec:samples}

The estimator above is defined for any set of occasions and weights. The main specification uses first adoptions
with $w_k = 1$; \cref{sec:results} reports the alternatives defined here, and a rule for calling the results robust
that we fixed before estimating them.

\paragraph{Samples.} \emph{P}: first adoptions (main). \emph{D}: all 24-hour decisions, including repeat decisions.
\emph{U}: all uploads. In D and U the author's earlier parents enter $\D_k$ with $\pi = 1$, so \cref{prop:pps}
applies unchanged. On a repeat decision the degree of a candidate includes the author's own earlier children of it,
which would load habit onto popularity. We therefore replace $d_{kl}$ by $\tilde d_{kl} = d_{kl} - o_{kl}$, where
$o_{kl}$ counts the children of $l$ by author $a_k$ created before $t_k$ (and similarly for adoptions in the last 30
days). On P, $o_{kl} \equiv 0$.

\paragraph{Weights.} With $N$ occasions by $A$ authors and $n_a$ occasions of author $a$, \emph{author weights}
$w_k = (N/A)/n_{a_k}$ give every author the same total weight and a mean weight of one. They estimate the behavior of
the typical author rather than of the typical choice.

\paragraph{Author types.} Authors are classified into six types (major organizations, re-distribution accounts,
other organizations, and individuals with $\ge$100, 10--99 and $<$10 models; \cref{app:data}). Heterogeneity is
tested by $LR = 2\big(\sum_\tau \ell_\tau(\hat\theta_\tau) - \ell(\hat\theta)\big)$, asymptotically $\chi^2$ with
degrees of freedom equal to the number of additional parameters of the type-specific models; because the logit scale may differ across types, types are compared by
shares, not coefficients.

\paragraph{Robustness rule.} An alternative sample or weighting is called robust if its popularity share is within 10
percentage points of the main specification and its fit/popularity ratio within a factor of three.
, and add
% % Proofs for sections/4_method_revised.tex. Input after the other appendices when the revised Section 4 is used.
\section{Proofs for Section~\ref{sec:method}}
\label{app:proofs}

\paragraph{Proof of \cref{lem:size}.}
(i) Let $I_{kl} = \1[l \in \D_k]$. For $l \ne j_k$, $I_{kl} \sim \mathrm{Bern}(\pi_{kl})$; $I_{kj_k} = 1$ because
the chosen parent is added if it was not drawn. By linearity of expectation (independence is not needed),
$\E|\D_k| = \sum_{l \ne j_k} \pi_{kl} + 1 = \sum_{l} \pi_{kl} + (1 - \pi_{kj_k})$, which lies in $[S, S+1]$ when
$\sum_l \pi_{kl} = S$.

(ii) Each term $n_\delta \min\{1, c(\delta+1)\}$ is continuous and non-decreasing in $c$, linear with slope
$n_\delta(\delta+1) > 0$ for $c < 1/(\delta+1)$ and constant afterwards; hence $\Sigma$ is continuous,
non-decreasing, strictly increasing wherever at least one class is unsaturated, $\Sigma(0) = 0$ and $\Sigma(c) =
|\M_k|$ once every class is saturated (for $c \ge 1$). If $|\M_k| > S$, the intermediate-value theorem gives a
solution, and at any solution some class is unsaturated (otherwise $\Sigma = |\M_k| > S$), so $\Sigma$ is strictly
increasing there and the solution is unique. Since $1/(\delta+1)$ decreases in $\delta$, the saturated classes at
$c_k$ are always the $m$ largest degrees. Given $m$, $\Sigma(c) = \sum_{i \le m} n_{\delta_i} + c\,(n_0 + \sum_{i >
m} n_{\delta_i}(\delta_i + 1))$, and solving $\Sigma(c) = S$ gives the formula. The consistency conditions select
the true $m$; trying $m = 0, 1, \dots$ in turn finds it in one pass over the distinct degrees, and the replay keeps
the degree histogram up to date incrementally. Degree-0 models are drawn as a $\mathrm{Bin}(n_0, c_k)$ count followed
by rejection sampling. \hfill$\square$

\paragraph{Proof of \cref{prop:pps}.}
Fix $k$ and drop the index. Every $l \ne j$ enters $\D$ independently with probability $\pi_l$, and $j$ is always in
$\D$, so given that $j$ was chosen the probability of a set $\D \ni j$ is $q(\D \mid j) = \prod_{l \in \D \setminus
\{j\}} \pi_l \prod_{l \notin \D}(1 - \pi_l)$. For $j, l \in \D$ the two products differ only in the factors of $j$
and $l$, so $q(\D \mid l) = q(\D \mid j)\,\pi_j / \pi_l$. By Bayes' rule and \eqref{eq:logit},
\[
  P(j \mid \D) = \frac{P(j \mid \M)\, q(\D \mid j)}{\sum_{l \in \D} P(l \mid \M)\, q(\D \mid l)}
  = \frac{e^{V_j}\, q(\D \mid j)}{\sum_{l \in \D} e^{V_l}\, q(\D \mid j)\, \pi_j / \pi_l}
  = \frac{e^{V_j - \log \pi_j}}{\sum_{l \in \D} e^{V_l - \log \pi_l}}.
\]
Only $P(j \mid \M) \propto e^{V_j}$ is used, so the result holds for any $V$ (in particular \eqref{eq:fe}).
\hfill$\square$

\paragraph{Proof of \cref{cor:consistency} (sketch).}
By \cref{prop:pps}, $\ell_k(\theta) = \log P(j_k \mid \D_k; \theta)$ is the log of a correctly specified conditional
density of the observed $j_k$ given $\D_k$, so $\theta_0$ maximizes $\E[\ell_k(\theta)]$ (information inequality).
By \cref{lem:concave}, $\ell_k$ is concave, and $-\nabla^2 \E[\ell_k] = \E[\mathrm{Cov}_{p_k}(x_k)]$ is positive
definite at $\theta_0$, so the maximizer is unique. Group occasions by author: the author sums
$L_a(\theta) = \sum_{k: a_k = a} \ell_k(\theta)$ are independent with bounded size, so a law of large numbers holds
pointwise, and for concave objectives pointwise convergence to a uniquely maximized limit implies consistency
\citep[Thm.~2.7]{newey1994large}. Asymptotic normality follows from a mean-value expansion of the first-order
condition and the central limit theorem for the independent author scores $\nabla L_a(\theta_0)$, giving the
sandwich $H_0^{-1} \Omega_0 H_0^{-1}$ \citep[Thm.~3.1]{newey1994large}. \hfill$\square$

\paragraph{Proof of \cref{lem:concave}.}
For one occasion, $\ell_k = v_{kj_k} - \log \sum_l e^{v_{kl}}$ with $\partial v_{kl} / \partial \theta = x_{kl}$.
Then $\nabla \ell_k = x_{kj_k} - \sum_l p_{kl} x_{kl}$, and differentiating $\bar x_k = \sum_l p_{kl} x_{kl}$ with
$\partial p_{kl} / \partial \theta = p_{kl}(x_{kl} - \bar x_k)$ gives $\nabla^2 \ell_k = -\sum_l p_{kl}(x_{kl} -
\bar x_k)(x_{kl} - \bar x_k)^\top = -\mathrm{Cov}_{p_k}(x_k)$, a negative semi-definite matrix. Summing with
non-negative weights preserves this; if the weighted sum of covariances is positive definite at every $\theta$, the
Hessian is negative definite and the maximizer is unique. \hfill$\square$

\paragraph{Proof of \cref{lem:noise}.}
With $\hat\phi_h = \phi_h + e_h$ and $\E[e_h \mid \phi_h, z_h, w_h] = 0$, $\mathrm{Cov}(\phi_h, e_h) = 0$, so
$\mathrm{Var}(\hat\phi) = \mathrm{Var}(\phi) + \E[\sigma_h^2]$. The population projection of $\hat\phi$ on $(z, w)$
equals that of $\phi$, because $\mathrm{Cov}(z, e) = \mathrm{Cov}(w, e) = 0$; hence $\mathrm{Var}(b^\top z)$ is the
variance of the quality part of $\phi$ itself, and dividing it by $\mathrm{Var}(\hat\phi) - \bar\sigma^2$ estimates
its share of $\mathrm{Var}(\phi)$. Without the correction the denominator is too large and the quality share too
small. We use the bootstrap variance of each $\hat\phi_h$ for $\sigma_h^2$ and keep hubs with bootstrap standard error
below one. \hfill$\square$

\paragraph{The proportional split.}
$s^{\mathrm{ext}}_{\mathrm{qual}} = s_{\mathrm{qual}} + q\, s_{\mathrm{FE}}$ assumes that the quality and the
non-quality parts of the hub intercepts contribute to the likelihood gain in proportion to their variances. An exact
allocation would enter the quality part $\hat b^\top z_h$ as a separate Shapley player, which requires external
quality for every hub; coverage is limited to \extHubs{} hubs. The lower end ($s_{\mathrm{qual}}$, no fixed effects)
and the upper end ($\bar s_{\mathrm{qual}}$, the whole fixed-effect share counted as quality) do not use the split.

\paragraph{The recovery draws.}
In \eqref{eq:sir}, $\D_k$ is the realized candidate set, which was drawn with the real $j_k$ forced in. The
simulated choice is therefore a draw from the model-implied conditional law on the realized set; the forced
inclusion of one real parent among about 300 candidates has a negligible effect on the composition of the set.
The draws use the four-group decomposition of the earlier analysis, whose inertia group equals self-supply and habit
combined; the full model, and hence $V^{\mathrm{FULL}}$, is the same under four and five groups.
 next to the other appendices (it holds the proofs cited here).
% All labels used elsewhere in the paper (sec:method, eq:logit, sec:sampling, eq:pps, prop:pps, eq:corrected,
% sec:shapley, eq:shapley, sec:bounds, sec:validation) are kept.
%
% Theorem environments: acmart defines definition, lemma, proposition, corollary; assumption and remark are
% defined below if missing (no change to main.tex needed).
%
% Changes relative to 4_method.tex, by comment of the author:
%  1. 4.1 now defines the choice occasion, degree and the information set (Def. 1) before the random-utility
%     assumption (Assumption 1); IIA is stated, because Prop. 1 uses it.
%  2. Consistency / asymptotic normality is a separate corollary (Cor. 1) after the sampling proposition.
%  3. 4.3 is about the estimator: weighted sampled likelihood, gradient and Hessian, concavity (Lemma 3), the
%     estimation algorithm (Algorithm 1), the clustered sandwich and the multiplier bootstrap. Hardware is one
%     sentence.
%  4. 4.4 starts from why coefficients cannot answer the question (units, groups, within-set variation, order
%     dependence with the real numbers) and defines the likelihood-gain game (Def. 3) before the Shapley value.
%  5. 4.5 defines the hub set formally (Def. 4) and states the noise-corrected split as a lemma.
%  6. 4.6 writes the three generating utilities, the simulated choice and the recovery criterion as equations.
%  New 4.7: samples, weights and author types (definitions only; results in Section 5, Table tab:samples).
%  Five groups (self-supply and habit replace inertia) and 2^5 = 32 fits throughout.
%
% Section 3 needs two matching edits (not done): (a) "the 500 most adopted parents" -> refer to Def. 4 (the hubs are
% ranked by degree d_l(t_h) over all children, not by first adoptions); (b) the inertia group -> self-supply
% (same_author) and habit (author_reuse, author_family, is_last).
% =====================================================================================================================
\makeatletter
\@ifundefined{assumption}{\newtheorem{assumption}{Assumption}}{}
\@ifundefined{remark}{\theoremstyle{acmdefinition}\newtheorem{remark}[theorem]{Remark}\theoremstyle{acmplain}}{}
\makeatother

\section{Model and Estimation}
\label{sec:method}

\Cref{sec:data} showed that adoption is highly concentrated and that popularity, general quality and demand match
could each produce this concentration. To separate them we need a model in which all three can act at once and be
compared on a common scale. This section builds that model (\cref{sec:choice}), makes it computable on markets of
hundreds of thousands of alternatives (\cref{sec:sampling,sec:estimation}), turns its estimates into a
decomposition by signal group (\cref{sec:shapley}), bounds what unmeasured quality could add (\cref{sec:bounds}),
and checks that the decomposition recovers known mechanisms (\cref{sec:validation}).

% ---------------------------------------------------------------------------------------------------------------------
\subsection{Base-model selection as a discrete choice}
\label{sec:choice}

\begin{definition}[Choice occasion]
\label{def:occasion}
A \emph{choice occasion} $k$ is a first adoption (\cref{sec:data}), described by its author $a_k$, time $t_k$,
modality $m_k$ and chosen parent $j_k$. Its \emph{market} $\M_k$ is the set of models of modality $m_k$ created
before $t_k$ that $a_k$ has not used as a parent. The \emph{degree} of a model $l$ at time $t$ is
\[
  d_l(t) = \#\{c : \text{$c$ declares $l$ as its parent with type fine-tune, adapter or unspecified, and $c$ was
  created before $t$}\},
\]
and $d_{kl} = d_l(t_k)$. Every candidate $l \in \M_k$ carries a feature vector $x_{kl} \in \mathbb{R}^p$ computed
only from models, edges and metadata that exist before $t_k$.
\end{definition}

The features are the five groups $g \in \mathcal G = \{\text{pop}, \text{qual}, \text{match}, \text{self},
\text{habit}\}$ of \cref{sec:data}: popularity, general quality, demand match, self-supply (the candidate is the
author's own model) and habit (the candidate belongs to a family the author used before). We write $x_{kl} =
(x_{kl,g})_{g \in \mathcal G}$ and $\theta = (\theta_g)_{g \in \mathcal G}$ conformably. Popularity and quality
features vary over candidates and time; demand-match, self-supply and habit features also vary over authors.

\begin{assumption}[Random utility]
\label{ass:rum}
The author of occasion $k$ chooses $j_k = \arg\max_{l \in \M_k} U_{kl}$, where
\[
  U_{kl} = V_{kl} + \varepsilon_{kl}, \qquad V_{kl} = \theta^\top x_{kl},
\]
and the $\varepsilon_{kl}$ are i.i.d.\ standard Type-I extreme value, independent of $\{x_{kl}\}_{l \in \M_k}$.
\end{assumption}

\begin{lemma}[\citealp{mcfadden1974conditional}]
\label{lem:logit}
Under \cref{ass:rum},
\begin{equation}
  P(j \mid \M_k) = \frac{\exp(\theta^\top x_{kj})}{\sum_{l \in \M_k} \exp(\theta^\top x_{kl})}.
  \label{eq:logit}
\end{equation}
\end{lemma}

\begin{remark}[What the model does and does not say]
\label{rem:logit}
(i) \emph{Identification.} Only differences across candidates within an occasion enter \eqref{eq:logit}; a
variable that is constant across $l \in \M_k$ (author experience, organizational status) cancels. Author
heterogeneity therefore enters through author-specific features (demand match, self-supply, habit) and through
estimation by author type (\cref{sec:samples}).
(ii) \emph{Scale.} The error variance is normalized, so $\theta$ is measured relative to the dispersion of tastes;
comparisons across samples are made with likelihood shares (\cref{sec:shapley}), which are scale free, not with
coefficients.
(iii) \emph{Independence of irrelevant alternatives.} $P(j \mid \M_k) / P(l \mid \M_k) = \exp(V_{kj} - V_{kl})$
does not depend on the other alternatives. \Cref{prop:pps} uses this property.
(iv) \emph{Interpretation.} We read $\theta$ descriptively, as the weights that best summarize which signals
choices co-move with, not as causal effects (\cref{sec:discussion}).
\end{remark}

Two network-growth models are special cases. With $\theta = 0$, \eqref{eq:logit} is \emph{uniform attachment}:
every model in the market is equally likely. With only $x_{kl} = (\log \max(d_{kl},1),\ \1[d_{kl}>0])$ and
coefficients $(\alpha, \beta)$,
\[
  P(j \mid \M_k) \;\propto\;
  \begin{cases} e^{\beta}\, d_{kj}^{\alpha}, & d_{kj} > 0, \\ 1, & d_{kj} = 0, \end{cases}
\]
which is \emph{preferential attachment} with exponent $\alpha$; $\alpha = 1$ is the linear model of
\citet{barabasi1999emergence}. Following \citet{overgoor2019choosing}, \eqref{eq:logit} nests these mechanisms and
lets them be compared by likelihood together with quality and demand match.

% ---------------------------------------------------------------------------------------------------------------------
\subsection{Sampling the competitors that matter}
\label{sec:sampling}

The denominator of \eqref{eq:logit} has a median of \marketMedian{} terms, and the training data contain \nTrain{}
occasions, so the likelihood cannot be evaluated on full markets. We evaluate it on a sampled candidate set.

\begin{definition}[Degree-proportional Poisson design]
\label{def:pps}
For occasion $k$ and a target size $S$, every $l \in \M_k$ enters the candidate set $\D_k$ independently with
probability
\begin{equation}
  \pi_{kl} = \min\{1,\; c_k (d_{kl} + 1)\},
  \label{eq:pps}
\end{equation}
where the scale $c_k > 0$ solves $\sum_{l \in \M_k} \pi_{kl} = S$. The chosen parent $j_k$ is added to $\D_k$ if it
was not drawn.
\end{definition}

The inclusion probability is proportional to $d+1$ (the $+1$ keeps never-used models in the design), and $c_k$ is
the factor of proportionality. The condition on $c_k$ fixes the expected size of the set:

\begin{lemma}[Set size and the scale $c_k$]
\label{lem:size}
(i) $\E|\D_k| = \sum_{l \in \M_k} \pi_{kl} + (1 - \pi_{kj_k}) \in [S, S+1]$.
(ii) Let $n_\delta$ be the number of models in $\M_k$ with degree $\delta$, and let $\delta_1 > \delta_2 > \dots >
\delta_M > 0$ be the positive degrees. The function $\Sigma(c) = \sum_\delta n_\delta \min\{1, c(\delta+1)\}$ is
continuous and non-decreasing, strictly increasing wherever some degree class is unsaturated
($c(\delta+1) < 1$), and runs from $0$ to $|\M_k|$. If $|\M_k| > S$, $\Sigma(c) = S$ has a unique solution
\[
  c_k = \frac{S - \sum_{i \le m} n_{\delta_i}}{n_0 + \sum_{i > m} n_{\delta_i}(\delta_i + 1)},
\]
where $m \in \{0, \dots, M\}$ is the number of saturated classes: the unique $m$ with $c_k(\delta_i + 1) \ge 1$ for
$i \le m$, $c_k(\delta_i + 1) < 1$ for $i > m$, and $c_k \le 1$.
\end{lemma}

\noindent The proof is in \cref{app:proofs}. With $S = 300$, on average \piOneMean{} candidates per training
occasion have $\pi = 1$; these are the incumbents a developer is most likely to weigh, and \piOneChosen\% of chosen
parents are among them. The rest of each set is a probability-proportional-to-size sample of the long tail.

\begin{proposition}[Sampling correction]
\label{prop:pps}
Under \cref{ass:rum} and \cref{def:pps},
\begin{equation}
  P(j \mid \D_k) = \frac{\exp(\theta^\top x_{kj} - \log \pi_{kj})}{\sum_{l \in \D_k} \exp(\theta^\top x_{kl} - \log
  \pi_{kl})}, \qquad j \in \D_k.
  \label{eq:corrected}
\end{equation}
\end{proposition}

\noindent Given that $j$ was chosen, the probability of drawing a set $\D \ni j$ is $q(\D \mid j) = \prod_{l \in \D
\setminus \{j\}} \pi_l \prod_{l \notin \D} (1 - \pi_l)$, so $q(\D \mid j)/q(\D \mid l) = \pi_l / \pi_j$ for any $j,
l \in \D$; Bayes' rule and \eqref{eq:logit} give \eqref{eq:corrected} (\cref{app:proofs};
\citealp{mcfadden1978modelling}). The argument uses only the logit form of $P(j \mid \M_k)$, so it holds for any
specification of $V_{kl}$, including the fixed-effect model of \cref{sec:bounds}.

\begin{corollary}[Consistency and asymptotic normality]
\label{cor:consistency}
Let $\hat\theta$ maximize $\ell(\theta) = \sum_k \log P(j_k \mid \D_k; \theta)$ from \eqref{eq:corrected}. Suppose
occasions are independent across authors with bounded numbers per author, the parameter space is compact with the
true $\theta_0$ in its interior, $\E\|x_{kl}\|^2 < \infty$, and $H_0 = \E[\mathrm{Cov}_{p_k}(x_k)]$ is positive
definite, where $\mathrm{Cov}_{p_k}$ is the covariance of $x_{kl}$ over $l \in \D_k$ under the probabilities
\eqref{eq:corrected}. Then, as the number of authors $A \to \infty$, $\hat\theta \to \theta_0$ in probability and
$\sqrt{A}(\hat\theta - \theta_0) \to N(0, H_0^{-1} \Omega_0 H_0^{-1})$, with $\Omega_0$ the variance of an author's
summed score.
\end{corollary}

\noindent By \cref{prop:pps}, $\ell$ is a correctly specified conditional likelihood of the observed $(j_k, \D_k)$,
and it is concave (\cref{lem:concave} below); the result follows from the standard argument for concave M-estimators
\citep{newey1994large}. Two consequences of \eqref{eq:corrected} matter for what follows.

\begin{remark}[Uniform attachment as the baseline, and why sample by degree]
\label{rem:pps}
(i) At $\theta = 0$, \eqref{eq:corrected} gives each sampled candidate the weight $1/\pi_{kl}$, and $\sum_{l \in
\D_k} 1/\pi_{kl}$ is a Horvitz--Thompson estimate of $|\M_k|$ \citep{horvitz1952generalization}; uniform attachment
on the sample therefore corresponds to uniform attachment on the full market. It is the baseline of every likelihood
gain below.
(ii) \Cref{cor:consistency} holds for any design with $\pi > 0$, but the information matrix $H_0 = \E[\mathrm{Cov}_{p_k}
(x_k)]$ depends on the design: a candidate set contributes information only through the variation of $x$ among the
candidates that have non-negligible choice probability. A uniform sample of 300 from $3 \times 10^5$ models almost
never contains the few models that receive most adoptions, so the chosen parent is compared mostly with never-used
models. Under a uniform design the median clustered standard error is \seRatioUnif{} times larger, and the
estimated weights themselves change (\cref{sec:robustq}).
\end{remark}

% ---------------------------------------------------------------------------------------------------------------------
\subsection{Estimation and inference}
\label{sec:estimation}

\Cref{cor:consistency} defines the estimator; we now state how it is computed and how its uncertainty is measured.
Write $v_{kl}(\theta) = \theta^\top x_{kl} - \log \pi_{kl}$, $p_{kl}(\theta) = e^{v_{kl}} / \sum_{m \in \D_k}
e^{v_{km}}$ and $\bar x_k(\theta) = \sum_{l \in \D_k} p_{kl} x_{kl}$. With occasion weights $w_k$ ($w_k = 1$ in the
main specification; other weights in \cref{sec:samples} and the bootstrap below), the sampled log-likelihood is
\begin{equation}
  \ell(\theta) = \sum_k w_k \Big[ v_{kj_k}(\theta) - \log \sum_{l \in \D_k} e^{v_{kl}(\theta)} \Big].
  \label{eq:loglik}
\end{equation}

\begin{lemma}[Concavity]
\label{lem:concave}
$\nabla \ell(\theta) = \sum_k w_k (x_{kj_k} - \bar x_k)$ and $\nabla^2 \ell(\theta) = -\sum_k w_k
\mathrm{Cov}_{p_k}(x_k)$. Hence $\ell$ is concave, and strictly concave, with a unique maximizer, whenever
$\sum_k w_k \mathrm{Cov}_{p_k}(x_k)$ is positive definite.
\end{lemma}

\noindent The offset $-\log \pi$ does not depend on $\theta$, so it changes neither the gradient's form nor the
concavity. The first-order condition equates the weighted mean features of the chosen parents with their mean under
the model. Concavity has a practical consequence: every refit in the decomposition and the bootstrap converges to
the same optimum from any starting point, so warm starts change only computing time.

\begin{center}
\fbox{\begin{minipage}{0.94\linewidth}
\textbf{Algorithm 1 (estimation).}
\begin{enumerate}
\item \emph{Replay.} Process the network in time order; at each occasion $k$ update $d_l(t_k)$, solve $c_k$
  (\cref{lem:size}), draw $\D_k$ by \cref{def:pps}, and record $(x_{kl}, \log \pi_{kl})_{l \in \D_k}$.
\item \emph{Fit.} Maximize \eqref{eq:loglik} by L-BFGS with the analytic gradient of \cref{lem:concave}; compute
  the Hessian at $\hat\theta$.
\item \emph{Subsets.} For every $\mathcal S \subseteq \mathcal G$, maximize \eqref{eq:loglik} with $\theta_g = 0$ for
  $g \notin \mathcal S$, warm-started from the fit of $\mathcal S$ minus one group ($2^{|\mathcal G|} = \nShapFits{}$
  fits; \cref{sec:shapley}).
\item \emph{Bootstrap.} Repeat steps 2--3 $B = 50$ times with weights $w_k \xi_{a_k}$, $\xi_a \sim \mathrm{Exp}(1)$
  i.i.d.\ over authors, warm-started from the point estimates.
\end{enumerate}
\end{minipage}}
\end{center}

\noindent The training data have \nRowsTrain{} million (occasion, candidate) rows; the full model fits in about a
minute on one GPU.\footnote{Candidate sets are grouped into buckets of similar size and kept on the GPU in single
precision; the likelihood, gradient and Hessian are accumulated in double precision.}

\paragraph{Inference.}
Occasions of one author are dependent (the same project, the same tastes, batch uploads): the \nFirst{} first
adoptions come from \nFirstAuthors{} authors, and the median clustered standard error is \seClusterInfl{} times
the one that treats occasions as independent. We therefore use the cluster-robust sandwich
\[
  \widehat{\mathrm{Var}}(\hat\theta) = \hat H^{-1} \Big( \textstyle\sum_a \hat s_a \hat s_a^\top \Big) \hat H^{-1},
  \qquad \hat H = -\nabla^2 \ell(\hat\theta), \quad \hat s_a = \sum_{k: a_k = a} w_k (x_{kj_k} - \bar x_k(\hat\theta)),
\]
\citep{cameron2015practitioner}. For derived quantities, the Shapley shares of \cref{sec:shapley} in particular, we
use the multiplier bootstrap of step 4: replicate $b$ maximizes $\sum_k \xi^{(b)}_{a_k} w_k \ell_k(\theta)$; the
shares are recomputed in each replicate; and intervals are the 2.5\% and 97.5\% percentiles. Normalized Exp(1)
weights are Dirichlet(1,\dots,1), the Bayesian bootstrap of \citet{rubin1981bayesian}; reweighting whole authors
preserves within-author dependence, and the weighted bootstrap is valid for M-estimators of this kind
\citep{ma2005robust}.

\paragraph{Out-of-sample evaluation.}
Let $K_h$ be the holdout occasions (the last 10\% in time) and $P_0$ uniform attachment ($\theta = 0$ in
\eqref{eq:corrected}). We report the mean holdout log-likelihood gain
\[
  \Delta_h = \frac{1}{|K_h|} \sum_{k \in K_h} \big[ \log P(j_k \mid \D_k; \hat\theta) - \log P_0(j_k \mid \D_k)
  \big],
\]
so that $e^{\Delta_h}$ is the geometric-mean factor by which the model raises the probability of the true parent
over uniform attachment ($\Delta_h = \holdGain$, a factor of \holdGainTimes), and the top-$k$ rate, the share of
holdout occasions in which $j_k$ ranks among the $k$ candidates with the highest $\hat V_{kl}$ (ranking by
utility, without the offset; ties broken at random).

% ---------------------------------------------------------------------------------------------------------------------
\subsection{From coefficients to a decomposition}
\label{sec:shapley}

$\hat\theta$ has 26 components, but the question is about groups of signals: how much of developers' choices do
popularity signals account for, and how much do quality and demand match? The coefficients cannot answer it
directly, for four reasons.
(i) Features have different units (indicators, logarithms of counts, similarities), so coefficients are not
comparable across features.
(ii) Each group contains several features, and the question concerns the group.
(iii) What a group contributes to explaining a choice depends on how much its \emph{index}
$I_{g,kl} = \theta_g^\top x_{kl,g}$ varies across the candidates of $\D_k$, not on the size of $\theta_g$ alone;
a large coefficient on a feature that rarely varies within a set explains little.
(iv) The groups are correlated: popular models are also well documented and match common tasks. The fit gained by
adding a group therefore depends on which groups are already in the model. Demand match accounts for \matchFirst\%
of the full model's likelihood gain when it enters first and for \matchLast\% when it enters last.
We therefore measure each group by its contribution to the model's fit and allocate that contribution with a rule
that does not depend on the order of entry.

\begin{definition}[Likelihood-gain game]
\label{def:game}
For $\mathcal S \subseteq \mathcal G$, let $\hat\theta_{\mathcal S}$ maximize \eqref{eq:loglik} subject to
$\theta_g = 0$ for $g \notin \mathcal S$, and let
\[
  G(\mathcal S) = \ell(\hat\theta_{\mathcal S}) - \ell(0),
\]
the log-likelihood gain over uniform attachment of the model with the groups in $\mathcal S$. The pair
$(\mathcal G, G)$ is a cooperative game with $G(\emptyset) = 0$. Its holdout version $G_h(\mathcal S)$ evaluates the
same $\hat\theta_{\mathcal S}$ by the mean holdout log-likelihood gain.
\end{definition}

Because the models are nested and fitted by maximum likelihood, $G$ is monotone: $\mathcal S \subseteq \mathcal T$
implies $G(\mathcal S) \le G(\mathcal T)$. We allocate $G(\mathcal G)$ by the Shapley value
\citep{shapley1953value},
\begin{equation}
  \phi_g = \sum_{\mathcal S \subseteq \mathcal G \setminus \{g\}}
  \frac{|\mathcal S|!\,(|\mathcal G|-|\mathcal S|-1)!}{|\mathcal G|!}
  \big[G(\mathcal S \cup \{g\}) - G(\mathcal S)\big],
  \qquad s_g = \frac{\phi_g}{G(\mathcal G)},
  \label{eq:shapley}
\end{equation}
the average marginal gain of group $g$ over all orders of entry. It is the only allocation that is efficient
($\sum_g \phi_g = G(\mathcal G)$), symmetric, assigns zero to a group that never adds fit, and is additive in $G$; by
monotonicity, every $s_g \in [0,1]$. Shapley allocations of fit are standard for correlated regressors
\citep{lipovetsky2001analysis,huettner2012axiomatic}. To report coarser groups, such as \emph{fit} = quality $\cup$
demand match, we apply \eqref{eq:shapley} to the game on blocks $v(\mathcal T) = G(\cup \mathcal T)$, whose values
are among the \nShapFits{} fits already computed. Shares are reported in-sample and on the holdout; on the holdout a
group gains only by predicting future choices, and $G_h$ need not be monotone.

% ---------------------------------------------------------------------------------------------------------------------
\subsection{Bounding general quality}
\label{sec:bounds}

Benchmark scores cover few models, so a small quality share could reflect missing measurement rather than
indifference to quality. We bound it from above with model-specific intercepts for the most used models.

\begin{definition}[Hubs]
\label{def:hubs}
Let $t_h$ be the train--holdout cut (\tHold). The \emph{hub set} $\mathcal H$ is the 500 models with the highest
degree $d_l(t_h)$ over the whole Hub (all modalities; ties broken by model index). $\mathcal H$ uses only training
data and does not depend on the sampling design.
\end{definition}

The fixed-effect model adds an intercept for each hub,
\begin{equation}
  V^{\mathrm{FE}}_{kl} = \theta^\top x_{kl} + \sum_{h \in \mathcal H} \phi_h \1[l = h],
  \label{eq:fe}
\end{equation}
estimated by \eqref{eq:corrected} with the same offset (\cref{prop:pps} holds for any $V$) and a small ridge penalty
$\frac{\lambda}{2}\|\phi\|^2$, $\lambda = 10^{-4}$. Non-hubs have $\phi = 0$, so $\phi_h$ is hub $h$'s
time-invariant attractiveness relative to the rest of the market. It absorbs every fixed attribute of the hub:
unmeasured quality, but also brand, tutorials and default status. Entering the fixed effects as an additional player
in \cref{def:game} gives a share $s_{\mathrm{FE}}$, and
\[
  \bar s_{\mathrm{qual}} = s_{\mathrm{qual}} + s_{\mathrm{FE}}
\]
is the largest share that any time-invariant hub attribute could contribute to quality (\feShare\% of the gain
in-sample).

To see how much of the fixed effects is quality, we project the estimated intercepts on external quality measures
$z_h$ (Open LLM Leaderboard v2 average; arena ratings) with controls $w_h$ (parameter count, re-distributor accounts),
$\hat\phi_h = b_0 + b^\top z_h + \gamma^\top w_h + u_h$, over the \extHubs{} hubs that are covered and precisely
estimated.

\begin{lemma}[Noise-corrected quality share]
\label{lem:noise}
Suppose $\hat\phi_h = \phi_h + e_h$ with $\E[e_h \mid \phi_h, z_h, w_h] = 0$ and $\mathrm{Var}(e_h) = \sigma_h^2$.
Then $\mathrm{Var}(\hat\phi) = \mathrm{Var}(\phi) + \bar\sigma^2$, the projection coefficients are unaffected by $e$,
and
\[
  q = \frac{\mathrm{Var}(\hat b^\top z_h)}{\mathrm{Var}(\hat\phi_h) - \bar\sigma^2}
\]
estimates the share of $\mathrm{Var}(\phi)$ explained by external quality, with $\sigma_h^2$ the bootstrap variance
of $\hat\phi_h$.
\end{lemma}

\noindent We split the fixed-effect share proportionally, $s^{\mathrm{ext}}_{\mathrm{qual}} = s_{\mathrm{qual}} + q\,
s_{\mathrm{FE}}$ ($q = \extQ\%$). Shapley shares are not linear in a split of a player, so this split is an
approximation; the two end points, $s_{\mathrm{qual}}$ and $\bar s_{\mathrm{qual}}$, do not depend on it
(\cref{app:proofs}).

% ---------------------------------------------------------------------------------------------------------------------
\subsection{Does the decomposition recover known mechanisms?}
\label{sec:validation}

The decomposition attributes most of the likelihood gain to popularity (\cref{sec:results}). Because the groups are
correlated, this could in principle be a property of the method rather than of the data. We test the method on
choices whose generating mechanism is known by construction.

For a generating utility $V^{(r)}$, keep every real candidate set $\D_k$ with its real features and inclusion
probabilities, and replace the observed parent $j_k$ by a draw
\begin{equation}
  \tilde j^{(r)}_k \sim P^{(r)}(l \mid \D_k) = \frac{\exp(V^{(r)}_{kl} - \log \pi_{kl})}{\sum_{m \in \D_k}
  \exp(V^{(r)}_{km} - \log \pi_{km})}, \qquad l \in \D_k,
  \label{eq:sir}
\end{equation}
the model-implied choice on the sampled set (\cref{prop:pps}). We use three utilities built from the estimates of
\cref{sec:estimation,sec:bounds}:
\[
  V^{\mathrm{SF}}_{kl} = \hat\theta_{\mathrm{pop}}^\top x_{kl,\mathrm{pop}}, \qquad
  V^{\mathrm{QF}}_{kl} = \hat\theta_{\mathrm{qual}}^\top x_{kl,\mathrm{qual}} + \hat\theta_{\mathrm{match}}^\top
  x_{kl,\mathrm{match}} + \hat q_l, \qquad
  V^{\mathrm{FULL}}_{kl} = \hat V^{\mathrm{FE}}_{kl},
\]
where $\hat q_l = \hat b^\top z_l$ is the part of the hub intercept explained by external quality ($\hat q_l = 0$ for
non-hubs). SF lets choices depend only on popularity, QF only on quality and demand match, and FULL on the estimated
model. Running steps 2--3 of Algorithm 1 on $\{(\tilde j^{(r)}_k, \D_k)\}_k$ gives shares $\hat s(r)$. If the
decomposition separates the mechanisms, $\hat s_{\mathrm{pop}}(\mathrm{SF})$ and $\hat s_{\mathrm{fit}}(\mathrm{QF})$
should be close to one, and $\hat s(\mathrm{FULL})$ should equal the shares estimated on the real choices; $1 -
\hat s_{\mathrm{pop}}(\mathrm{SF})$ and $1 - \hat s_{\mathrm{fit}}(\mathrm{QF})$ measure how much fit leaks to a
correlated group.

The decomposition attributes \recSfPop\% of the gain to popularity under SF and \recQfFit\% to fit under QF, and it
returns \recFullPop\%, \recFullMatch\% and \recFullInertia\% (popularity, demand match, self-supply plus habit) under
FULL, the shares estimated on the real choices.\todo{Recovery is run with four groups; rerun with five or keep the
note that self-supply + habit = the former inertia group.} When choices are generated by quality and demand match
alone, the decomposition does not manufacture a popularity share. As a check of the whole pipeline, we also grow
synthetic networks in which a new node attaches to $l$ with weight $d_l^{\alpha}$ if $d_l > 0$ and $e^{\beta_0}$
otherwise ($\alpha = 1$, $\beta_0 = -1$), replay them with the same code, and recover $\hat\alpha$ within three
standard errors of one under all three sampling designs (\cref{app:validation}).

% ---------------------------------------------------------------------------------------------------------------------
\subsection{Samples, weights and author types}
\label{sec:samples}

The estimator above is defined for any set of occasions and weights. The main specification uses first adoptions
with $w_k = 1$; \cref{sec:results} reports the alternatives defined here, and a rule for calling the results robust
that we fixed before estimating them.

\paragraph{Samples.} \emph{P}: first adoptions (main). \emph{D}: all 24-hour decisions, including repeat decisions.
\emph{U}: all uploads. In D and U the author's earlier parents enter $\D_k$ with $\pi = 1$, so \cref{prop:pps}
applies unchanged. On a repeat decision the degree of a candidate includes the author's own earlier children of it,
which would load habit onto popularity. We therefore replace $d_{kl}$ by $\tilde d_{kl} = d_{kl} - o_{kl}$, where
$o_{kl}$ counts the children of $l$ by author $a_k$ created before $t_k$ (and similarly for adoptions in the last 30
days). On P, $o_{kl} \equiv 0$.

\paragraph{Weights.} With $N$ occasions by $A$ authors and $n_a$ occasions of author $a$, \emph{author weights}
$w_k = (N/A)/n_{a_k}$ give every author the same total weight and a mean weight of one. They estimate the behavior of
the typical author rather than of the typical choice.

\paragraph{Author types.} Authors are classified into six types (major organizations, re-distribution accounts,
other organizations, and individuals with $\ge$100, 10--99 and $<$10 models; \cref{app:data}). Heterogeneity is
tested by $LR = 2\big(\sum_\tau \ell_\tau(\hat\theta_\tau) - \ell(\hat\theta)\big)$, asymptotically $\chi^2$ with
degrees of freedom equal to the number of additional parameters of the type-specific models; because the logit scale may differ across types, types are compared by
shares, not coefficients.

\paragraph{Robustness rule.} An alternative sample or weighting is called robust if its popularity share is within 10
percentage points of the main specification and its fit/popularity ratio within a factor of three.
, and add
% % Proofs for sections/4_method_revised.tex. Input after the other appendices when the revised Section 4 is used.
\section{Proofs for Section~\ref{sec:method}}
\label{app:proofs}

\paragraph{Proof of \cref{lem:size}.}
(i) Let $I_{kl} = \1[l \in \D_k]$. For $l \ne j_k$, $I_{kl} \sim \mathrm{Bern}(\pi_{kl})$; $I_{kj_k} = 1$ because
the chosen parent is added if it was not drawn. By linearity of expectation (independence is not needed),
$\E|\D_k| = \sum_{l \ne j_k} \pi_{kl} + 1 = \sum_{l} \pi_{kl} + (1 - \pi_{kj_k})$, which lies in $[S, S+1]$ when
$\sum_l \pi_{kl} = S$.

(ii) Each term $n_\delta \min\{1, c(\delta+1)\}$ is continuous and non-decreasing in $c$, linear with slope
$n_\delta(\delta+1) > 0$ for $c < 1/(\delta+1)$ and constant afterwards; hence $\Sigma$ is continuous,
non-decreasing, strictly increasing wherever at least one class is unsaturated, $\Sigma(0) = 0$ and $\Sigma(c) =
|\M_k|$ once every class is saturated (for $c \ge 1$). If $|\M_k| > S$, the intermediate-value theorem gives a
solution, and at any solution some class is unsaturated (otherwise $\Sigma = |\M_k| > S$), so $\Sigma$ is strictly
increasing there and the solution is unique. Since $1/(\delta+1)$ decreases in $\delta$, the saturated classes at
$c_k$ are always the $m$ largest degrees. Given $m$, $\Sigma(c) = \sum_{i \le m} n_{\delta_i} + c\,(n_0 + \sum_{i >
m} n_{\delta_i}(\delta_i + 1))$, and solving $\Sigma(c) = S$ gives the formula. The consistency conditions select
the true $m$; trying $m = 0, 1, \dots$ in turn finds it in one pass over the distinct degrees, and the replay keeps
the degree histogram up to date incrementally. Degree-0 models are drawn as a $\mathrm{Bin}(n_0, c_k)$ count followed
by rejection sampling. \hfill$\square$

\paragraph{Proof of \cref{prop:pps}.}
Fix $k$ and drop the index. Every $l \ne j$ enters $\D$ independently with probability $\pi_l$, and $j$ is always in
$\D$, so given that $j$ was chosen the probability of a set $\D \ni j$ is $q(\D \mid j) = \prod_{l \in \D \setminus
\{j\}} \pi_l \prod_{l \notin \D}(1 - \pi_l)$. For $j, l \in \D$ the two products differ only in the factors of $j$
and $l$, so $q(\D \mid l) = q(\D \mid j)\,\pi_j / \pi_l$. By Bayes' rule and \eqref{eq:logit},
\[
  P(j \mid \D) = \frac{P(j \mid \M)\, q(\D \mid j)}{\sum_{l \in \D} P(l \mid \M)\, q(\D \mid l)}
  = \frac{e^{V_j}\, q(\D \mid j)}{\sum_{l \in \D} e^{V_l}\, q(\D \mid j)\, \pi_j / \pi_l}
  = \frac{e^{V_j - \log \pi_j}}{\sum_{l \in \D} e^{V_l - \log \pi_l}}.
\]
Only $P(j \mid \M) \propto e^{V_j}$ is used, so the result holds for any $V$ (in particular \eqref{eq:fe}).
\hfill$\square$

\paragraph{Proof of \cref{cor:consistency} (sketch).}
By \cref{prop:pps}, $\ell_k(\theta) = \log P(j_k \mid \D_k; \theta)$ is the log of a correctly specified conditional
density of the observed $j_k$ given $\D_k$, so $\theta_0$ maximizes $\E[\ell_k(\theta)]$ (information inequality).
By \cref{lem:concave}, $\ell_k$ is concave, and $-\nabla^2 \E[\ell_k] = \E[\mathrm{Cov}_{p_k}(x_k)]$ is positive
definite at $\theta_0$, so the maximizer is unique. Group occasions by author: the author sums
$L_a(\theta) = \sum_{k: a_k = a} \ell_k(\theta)$ are independent with bounded size, so a law of large numbers holds
pointwise, and for concave objectives pointwise convergence to a uniquely maximized limit implies consistency
\citep[Thm.~2.7]{newey1994large}. Asymptotic normality follows from a mean-value expansion of the first-order
condition and the central limit theorem for the independent author scores $\nabla L_a(\theta_0)$, giving the
sandwich $H_0^{-1} \Omega_0 H_0^{-1}$ \citep[Thm.~3.1]{newey1994large}. \hfill$\square$

\paragraph{Proof of \cref{lem:concave}.}
For one occasion, $\ell_k = v_{kj_k} - \log \sum_l e^{v_{kl}}$ with $\partial v_{kl} / \partial \theta = x_{kl}$.
Then $\nabla \ell_k = x_{kj_k} - \sum_l p_{kl} x_{kl}$, and differentiating $\bar x_k = \sum_l p_{kl} x_{kl}$ with
$\partial p_{kl} / \partial \theta = p_{kl}(x_{kl} - \bar x_k)$ gives $\nabla^2 \ell_k = -\sum_l p_{kl}(x_{kl} -
\bar x_k)(x_{kl} - \bar x_k)^\top = -\mathrm{Cov}_{p_k}(x_k)$, a negative semi-definite matrix. Summing with
non-negative weights preserves this; if the weighted sum of covariances is positive definite at every $\theta$, the
Hessian is negative definite and the maximizer is unique. \hfill$\square$

\paragraph{Proof of \cref{lem:noise}.}
With $\hat\phi_h = \phi_h + e_h$ and $\E[e_h \mid \phi_h, z_h, w_h] = 0$, $\mathrm{Cov}(\phi_h, e_h) = 0$, so
$\mathrm{Var}(\hat\phi) = \mathrm{Var}(\phi) + \E[\sigma_h^2]$. The population projection of $\hat\phi$ on $(z, w)$
equals that of $\phi$, because $\mathrm{Cov}(z, e) = \mathrm{Cov}(w, e) = 0$; hence $\mathrm{Var}(b^\top z)$ is the
variance of the quality part of $\phi$ itself, and dividing it by $\mathrm{Var}(\hat\phi) - \bar\sigma^2$ estimates
its share of $\mathrm{Var}(\phi)$. Without the correction the denominator is too large and the quality share too
small. We use the bootstrap variance of each $\hat\phi_h$ for $\sigma_h^2$ and keep hubs with bootstrap standard error
below one. \hfill$\square$

\paragraph{The proportional split.}
$s^{\mathrm{ext}}_{\mathrm{qual}} = s_{\mathrm{qual}} + q\, s_{\mathrm{FE}}$ assumes that the quality and the
non-quality parts of the hub intercepts contribute to the likelihood gain in proportion to their variances. An exact
allocation would enter the quality part $\hat b^\top z_h$ as a separate Shapley player, which requires external
quality for every hub; coverage is limited to \extHubs{} hubs. The lower end ($s_{\mathrm{qual}}$, no fixed effects)
and the upper end ($\bar s_{\mathrm{qual}}$, the whole fixed-effect share counted as quality) do not use the split.

\paragraph{The recovery draws.}
In \eqref{eq:sir}, $\D_k$ is the realized candidate set, which was drawn with the real $j_k$ forced in. The
simulated choice is therefore a draw from the model-implied conditional law on the realized set; the forced
inclusion of one real parent among about 300 candidates has a negligible effect on the composition of the set.
The draws use the four-group decomposition of the earlier analysis, whose inertia group equals self-supply and habit
combined; the full model, and hence $V^{\mathrm{FULL}}$, is the same under four and five groups.
 next to the other appendices (it holds the proofs cited here).
% All labels used elsewhere in the paper (sec:method, eq:logit, sec:sampling, eq:pps, prop:pps, eq:corrected,
% sec:shapley, eq:shapley, sec:bounds, sec:validation) are kept.
%
% Theorem environments: acmart defines definition, lemma, proposition, corollary; assumption and remark are
% defined below if missing (no change to main.tex needed).
%
% Changes relative to 4_method.tex, by comment of the author:
%  1. 4.1 now defines the choice occasion, degree and the information set (Def. 1) before the random-utility
%     assumption (Assumption 1); IIA is stated, because Prop. 1 uses it.
%  2. Consistency / asymptotic normality is a separate corollary (Cor. 1) after the sampling proposition.
%  3. 4.3 is about the estimator: weighted sampled likelihood, gradient and Hessian, concavity (Lemma 3), the
%     estimation algorithm (Algorithm 1), the clustered sandwich and the multiplier bootstrap. Hardware is one
%     sentence.
%  4. 4.4 starts from why coefficients cannot answer the question (units, groups, within-set variation, order
%     dependence with the real numbers) and defines the likelihood-gain game (Def. 3) before the Shapley value.
%  5. 4.5 defines the hub set formally (Def. 4) and states the noise-corrected split as a lemma.
%  6. 4.6 writes the three generating utilities, the simulated choice and the recovery criterion as equations.
%  New 4.7: samples, weights and author types (definitions only; results in Section 5, Table tab:samples).
%  Five groups (self-supply and habit replace inertia) and 2^5 = 32 fits throughout.
%
% Section 3 needs two matching edits (not done): (a) "the 500 most adopted parents" -> refer to Def. 4 (the hubs are
% ranked by degree d_l(t_h) over all children, not by first adoptions); (b) the inertia group -> self-supply
% (same_author) and habit (author_reuse, author_family, is_last).
% =====================================================================================================================
\makeatletter
\@ifundefined{assumption}{\newtheorem{assumption}{Assumption}}{}
\@ifundefined{remark}{\theoremstyle{acmdefinition}\newtheorem{remark}[theorem]{Remark}\theoremstyle{acmplain}}{}
\makeatother

\section{Model and Estimation}
\label{sec:method}

\Cref{sec:data} showed that adoption is highly concentrated and that popularity, general quality and demand match
could each produce this concentration. To separate them we need a model in which all three can act at once and be
compared on a common scale. This section builds that model (\cref{sec:choice}), makes it computable on markets of
hundreds of thousands of alternatives (\cref{sec:sampling,sec:estimation}), turns its estimates into a
decomposition by signal group (\cref{sec:shapley}), bounds what unmeasured quality could add (\cref{sec:bounds}),
and checks that the decomposition recovers known mechanisms (\cref{sec:validation}).

% ---------------------------------------------------------------------------------------------------------------------
\subsection{Base-model selection as a discrete choice}
\label{sec:choice}

\begin{definition}[Choice occasion]
\label{def:occasion}
A \emph{choice occasion} $k$ is a first adoption (\cref{sec:data}), described by its author $a_k$, time $t_k$,
modality $m_k$ and chosen parent $j_k$. Its \emph{market} $\M_k$ is the set of models of modality $m_k$ created
before $t_k$ that $a_k$ has not used as a parent. The \emph{degree} of a model $l$ at time $t$ is
\[
  d_l(t) = \#\{c : \text{$c$ declares $l$ as its parent with type fine-tune, adapter or unspecified, and $c$ was
  created before $t$}\},
\]
and $d_{kl} = d_l(t_k)$. Every candidate $l \in \M_k$ carries a feature vector $x_{kl} \in \mathbb{R}^p$ computed
only from models, edges and metadata that exist before $t_k$.
\end{definition}

The features are the five groups $g \in \mathcal G = \{\text{pop}, \text{qual}, \text{match}, \text{self},
\text{habit}\}$ of \cref{sec:data}: popularity, general quality, demand match, self-supply (the candidate is the
author's own model) and habit (the candidate belongs to a family the author used before). We write $x_{kl} =
(x_{kl,g})_{g \in \mathcal G}$ and $\theta = (\theta_g)_{g \in \mathcal G}$ conformably. Popularity and quality
features vary over candidates and time; demand-match, self-supply and habit features also vary over authors.

\begin{assumption}[Random utility]
\label{ass:rum}
The author of occasion $k$ chooses $j_k = \arg\max_{l \in \M_k} U_{kl}$, where
\[
  U_{kl} = V_{kl} + \varepsilon_{kl}, \qquad V_{kl} = \theta^\top x_{kl},
\]
and the $\varepsilon_{kl}$ are i.i.d.\ standard Type-I extreme value, independent of $\{x_{kl}\}_{l \in \M_k}$.
\end{assumption}

\begin{lemma}[\citealp{mcfadden1974conditional}]
\label{lem:logit}
Under \cref{ass:rum},
\begin{equation}
  P(j \mid \M_k) = \frac{\exp(\theta^\top x_{kj})}{\sum_{l \in \M_k} \exp(\theta^\top x_{kl})}.
  \label{eq:logit}
\end{equation}
\end{lemma}

\begin{remark}[What the model does and does not say]
\label{rem:logit}
(i) \emph{Identification.} Only differences across candidates within an occasion enter \eqref{eq:logit}; a
variable that is constant across $l \in \M_k$ (author experience, organizational status) cancels. Author
heterogeneity therefore enters through author-specific features (demand match, self-supply, habit) and through
estimation by author type (\cref{sec:samples}).
(ii) \emph{Scale.} The error variance is normalized, so $\theta$ is measured relative to the dispersion of tastes;
comparisons across samples are made with likelihood shares (\cref{sec:shapley}), which are scale free, not with
coefficients.
(iii) \emph{Independence of irrelevant alternatives.} $P(j \mid \M_k) / P(l \mid \M_k) = \exp(V_{kj} - V_{kl})$
does not depend on the other alternatives. \Cref{prop:pps} uses this property.
(iv) \emph{Interpretation.} We read $\theta$ descriptively, as the weights that best summarize which signals
choices co-move with, not as causal effects (\cref{sec:discussion}).
\end{remark}

Two network-growth models are special cases. With $\theta = 0$, \eqref{eq:logit} is \emph{uniform attachment}:
every model in the market is equally likely. With only $x_{kl} = (\log \max(d_{kl},1),\ \1[d_{kl}>0])$ and
coefficients $(\alpha, \beta)$,
\[
  P(j \mid \M_k) \;\propto\;
  \begin{cases} e^{\beta}\, d_{kj}^{\alpha}, & d_{kj} > 0, \\ 1, & d_{kj} = 0, \end{cases}
\]
which is \emph{preferential attachment} with exponent $\alpha$; $\alpha = 1$ is the linear model of
\citet{barabasi1999emergence}. Following \citet{overgoor2019choosing}, \eqref{eq:logit} nests these mechanisms and
lets them be compared by likelihood together with quality and demand match.

% ---------------------------------------------------------------------------------------------------------------------
\subsection{Sampling the competitors that matter}
\label{sec:sampling}

The denominator of \eqref{eq:logit} has a median of \marketMedian{} terms, and the training data contain \nTrain{}
occasions, so the likelihood cannot be evaluated on full markets. We evaluate it on a sampled candidate set.

\begin{definition}[Degree-proportional Poisson design]
\label{def:pps}
For occasion $k$ and a target size $S$, every $l \in \M_k$ enters the candidate set $\D_k$ independently with
probability
\begin{equation}
  \pi_{kl} = \min\{1,\; c_k (d_{kl} + 1)\},
  \label{eq:pps}
\end{equation}
where the scale $c_k > 0$ solves $\sum_{l \in \M_k} \pi_{kl} = S$. The chosen parent $j_k$ is added to $\D_k$ if it
was not drawn.
\end{definition}

The inclusion probability is proportional to $d+1$ (the $+1$ keeps never-used models in the design), and $c_k$ is
the factor of proportionality. The condition on $c_k$ fixes the expected size of the set:

\begin{lemma}[Set size and the scale $c_k$]
\label{lem:size}
(i) $\E|\D_k| = \sum_{l \in \M_k} \pi_{kl} + (1 - \pi_{kj_k}) \in [S, S+1]$.
(ii) Let $n_\delta$ be the number of models in $\M_k$ with degree $\delta$, and let $\delta_1 > \delta_2 > \dots >
\delta_M > 0$ be the positive degrees. The function $\Sigma(c) = \sum_\delta n_\delta \min\{1, c(\delta+1)\}$ is
continuous and non-decreasing, strictly increasing wherever some degree class is unsaturated
($c(\delta+1) < 1$), and runs from $0$ to $|\M_k|$. If $|\M_k| > S$, $\Sigma(c) = S$ has a unique solution
\[
  c_k = \frac{S - \sum_{i \le m} n_{\delta_i}}{n_0 + \sum_{i > m} n_{\delta_i}(\delta_i + 1)},
\]
where $m \in \{0, \dots, M\}$ is the number of saturated classes: the unique $m$ with $c_k(\delta_i + 1) \ge 1$ for
$i \le m$, $c_k(\delta_i + 1) < 1$ for $i > m$, and $c_k \le 1$.
\end{lemma}

\noindent The proof is in \cref{app:proofs}. With $S = 300$, on average \piOneMean{} candidates per training
occasion have $\pi = 1$; these are the incumbents a developer is most likely to weigh, and \piOneChosen\% of chosen
parents are among them. The rest of each set is a probability-proportional-to-size sample of the long tail.

\begin{proposition}[Sampling correction]
\label{prop:pps}
Under \cref{ass:rum} and \cref{def:pps},
\begin{equation}
  P(j \mid \D_k) = \frac{\exp(\theta^\top x_{kj} - \log \pi_{kj})}{\sum_{l \in \D_k} \exp(\theta^\top x_{kl} - \log
  \pi_{kl})}, \qquad j \in \D_k.
  \label{eq:corrected}
\end{equation}
\end{proposition}

\noindent Given that $j$ was chosen, the probability of drawing a set $\D \ni j$ is $q(\D \mid j) = \prod_{l \in \D
\setminus \{j\}} \pi_l \prod_{l \notin \D} (1 - \pi_l)$, so $q(\D \mid j)/q(\D \mid l) = \pi_l / \pi_j$ for any $j,
l \in \D$; Bayes' rule and \eqref{eq:logit} give \eqref{eq:corrected} (\cref{app:proofs};
\citealp{mcfadden1978modelling}). The argument uses only the logit form of $P(j \mid \M_k)$, so it holds for any
specification of $V_{kl}$, including the fixed-effect model of \cref{sec:bounds}.

\begin{corollary}[Consistency and asymptotic normality]
\label{cor:consistency}
Let $\hat\theta$ maximize $\ell(\theta) = \sum_k \log P(j_k \mid \D_k; \theta)$ from \eqref{eq:corrected}. Suppose
occasions are independent across authors with bounded numbers per author, the parameter space is compact with the
true $\theta_0$ in its interior, $\E\|x_{kl}\|^2 < \infty$, and $H_0 = \E[\mathrm{Cov}_{p_k}(x_k)]$ is positive
definite, where $\mathrm{Cov}_{p_k}$ is the covariance of $x_{kl}$ over $l \in \D_k$ under the probabilities
\eqref{eq:corrected}. Then, as the number of authors $A \to \infty$, $\hat\theta \to \theta_0$ in probability and
$\sqrt{A}(\hat\theta - \theta_0) \to N(0, H_0^{-1} \Omega_0 H_0^{-1})$, with $\Omega_0$ the variance of an author's
summed score.
\end{corollary}

\noindent By \cref{prop:pps}, $\ell$ is a correctly specified conditional likelihood of the observed $(j_k, \D_k)$,
and it is concave (\cref{lem:concave} below); the result follows from the standard argument for concave M-estimators
\citep{newey1994large}. Two consequences of \eqref{eq:corrected} matter for what follows.

\begin{remark}[Uniform attachment as the baseline, and why sample by degree]
\label{rem:pps}
(i) At $\theta = 0$, \eqref{eq:corrected} gives each sampled candidate the weight $1/\pi_{kl}$, and $\sum_{l \in
\D_k} 1/\pi_{kl}$ is a Horvitz--Thompson estimate of $|\M_k|$ \citep{horvitz1952generalization}; uniform attachment
on the sample therefore corresponds to uniform attachment on the full market. It is the baseline of every likelihood
gain below.
(ii) \Cref{cor:consistency} holds for any design with $\pi > 0$, but the information matrix $H_0 = \E[\mathrm{Cov}_{p_k}
(x_k)]$ depends on the design: a candidate set contributes information only through the variation of $x$ among the
candidates that have non-negligible choice probability. A uniform sample of 300 from $3 \times 10^5$ models almost
never contains the few models that receive most adoptions, so the chosen parent is compared mostly with never-used
models. Under a uniform design the median clustered standard error is \seRatioUnif{} times larger, and the
estimated weights themselves change (\cref{sec:robustq}).
\end{remark}

% ---------------------------------------------------------------------------------------------------------------------
\subsection{Estimation and inference}
\label{sec:estimation}

\Cref{cor:consistency} defines the estimator; we now state how it is computed and how its uncertainty is measured.
Write $v_{kl}(\theta) = \theta^\top x_{kl} - \log \pi_{kl}$, $p_{kl}(\theta) = e^{v_{kl}} / \sum_{m \in \D_k}
e^{v_{km}}$ and $\bar x_k(\theta) = \sum_{l \in \D_k} p_{kl} x_{kl}$. With occasion weights $w_k$ ($w_k = 1$ in the
main specification; other weights in \cref{sec:samples} and the bootstrap below), the sampled log-likelihood is
\begin{equation}
  \ell(\theta) = \sum_k w_k \Big[ v_{kj_k}(\theta) - \log \sum_{l \in \D_k} e^{v_{kl}(\theta)} \Big].
  \label{eq:loglik}
\end{equation}

\begin{lemma}[Concavity]
\label{lem:concave}
$\nabla \ell(\theta) = \sum_k w_k (x_{kj_k} - \bar x_k)$ and $\nabla^2 \ell(\theta) = -\sum_k w_k
\mathrm{Cov}_{p_k}(x_k)$. Hence $\ell$ is concave, and strictly concave, with a unique maximizer, whenever
$\sum_k w_k \mathrm{Cov}_{p_k}(x_k)$ is positive definite.
\end{lemma}

\noindent The offset $-\log \pi$ does not depend on $\theta$, so it changes neither the gradient's form nor the
concavity. The first-order condition equates the weighted mean features of the chosen parents with their mean under
the model. Concavity has a practical consequence: every refit in the decomposition and the bootstrap converges to
the same optimum from any starting point, so warm starts change only computing time.

\begin{center}
\fbox{\begin{minipage}{0.94\linewidth}
\textbf{Algorithm 1 (estimation).}
\begin{enumerate}
\item \emph{Replay.} Process the network in time order; at each occasion $k$ update $d_l(t_k)$, solve $c_k$
  (\cref{lem:size}), draw $\D_k$ by \cref{def:pps}, and record $(x_{kl}, \log \pi_{kl})_{l \in \D_k}$.
\item \emph{Fit.} Maximize \eqref{eq:loglik} by L-BFGS with the analytic gradient of \cref{lem:concave}; compute
  the Hessian at $\hat\theta$.
\item \emph{Subsets.} For every $\mathcal S \subseteq \mathcal G$, maximize \eqref{eq:loglik} with $\theta_g = 0$ for
  $g \notin \mathcal S$, warm-started from the fit of $\mathcal S$ minus one group ($2^{|\mathcal G|} = \nShapFits{}$
  fits; \cref{sec:shapley}).
\item \emph{Bootstrap.} Repeat steps 2--3 $B = 50$ times with weights $w_k \xi_{a_k}$, $\xi_a \sim \mathrm{Exp}(1)$
  i.i.d.\ over authors, warm-started from the point estimates.
\end{enumerate}
\end{minipage}}
\end{center}

\noindent The training data have \nRowsTrain{} million (occasion, candidate) rows; the full model fits in about a
minute on one GPU.\footnote{Candidate sets are grouped into buckets of similar size and kept on the GPU in single
precision; the likelihood, gradient and Hessian are accumulated in double precision.}

\paragraph{Inference.}
Occasions of one author are dependent (the same project, the same tastes, batch uploads): the \nFirst{} first
adoptions come from \nFirstAuthors{} authors, and the median clustered standard error is \seClusterInfl{} times
the one that treats occasions as independent. We therefore use the cluster-robust sandwich
\[
  \widehat{\mathrm{Var}}(\hat\theta) = \hat H^{-1} \Big( \textstyle\sum_a \hat s_a \hat s_a^\top \Big) \hat H^{-1},
  \qquad \hat H = -\nabla^2 \ell(\hat\theta), \quad \hat s_a = \sum_{k: a_k = a} w_k (x_{kj_k} - \bar x_k(\hat\theta)),
\]
\citep{cameron2015practitioner}. For derived quantities, the Shapley shares of \cref{sec:shapley} in particular, we
use the multiplier bootstrap of step 4: replicate $b$ maximizes $\sum_k \xi^{(b)}_{a_k} w_k \ell_k(\theta)$; the
shares are recomputed in each replicate; and intervals are the 2.5\% and 97.5\% percentiles. Normalized Exp(1)
weights are Dirichlet(1,\dots,1), the Bayesian bootstrap of \citet{rubin1981bayesian}; reweighting whole authors
preserves within-author dependence, and the weighted bootstrap is valid for M-estimators of this kind
\citep{ma2005robust}.

\paragraph{Out-of-sample evaluation.}
Let $K_h$ be the holdout occasions (the last 10\% in time) and $P_0$ uniform attachment ($\theta = 0$ in
\eqref{eq:corrected}). We report the mean holdout log-likelihood gain
\[
  \Delta_h = \frac{1}{|K_h|} \sum_{k \in K_h} \big[ \log P(j_k \mid \D_k; \hat\theta) - \log P_0(j_k \mid \D_k)
  \big],
\]
so that $e^{\Delta_h}$ is the geometric-mean factor by which the model raises the probability of the true parent
over uniform attachment ($\Delta_h = \holdGain$, a factor of \holdGainTimes), and the top-$k$ rate, the share of
holdout occasions in which $j_k$ ranks among the $k$ candidates with the highest $\hat V_{kl}$ (ranking by
utility, without the offset; ties broken at random).

% ---------------------------------------------------------------------------------------------------------------------
\subsection{From coefficients to a decomposition}
\label{sec:shapley}

$\hat\theta$ has 26 components, but the question is about groups of signals: how much of developers' choices do
popularity signals account for, and how much do quality and demand match? The coefficients cannot answer it
directly, for four reasons.
(i) Features have different units (indicators, logarithms of counts, similarities), so coefficients are not
comparable across features.
(ii) Each group contains several features, and the question concerns the group.
(iii) What a group contributes to explaining a choice depends on how much its \emph{index}
$I_{g,kl} = \theta_g^\top x_{kl,g}$ varies across the candidates of $\D_k$, not on the size of $\theta_g$ alone;
a large coefficient on a feature that rarely varies within a set explains little.
(iv) The groups are correlated: popular models are also well documented and match common tasks. The fit gained by
adding a group therefore depends on which groups are already in the model. Demand match accounts for \matchFirst\%
of the full model's likelihood gain when it enters first and for \matchLast\% when it enters last.
We therefore measure each group by its contribution to the model's fit and allocate that contribution with a rule
that does not depend on the order of entry.

\begin{definition}[Likelihood-gain game]
\label{def:game}
For $\mathcal S \subseteq \mathcal G$, let $\hat\theta_{\mathcal S}$ maximize \eqref{eq:loglik} subject to
$\theta_g = 0$ for $g \notin \mathcal S$, and let
\[
  G(\mathcal S) = \ell(\hat\theta_{\mathcal S}) - \ell(0),
\]
the log-likelihood gain over uniform attachment of the model with the groups in $\mathcal S$. The pair
$(\mathcal G, G)$ is a cooperative game with $G(\emptyset) = 0$. Its holdout version $G_h(\mathcal S)$ evaluates the
same $\hat\theta_{\mathcal S}$ by the mean holdout log-likelihood gain.
\end{definition}

Because the models are nested and fitted by maximum likelihood, $G$ is monotone: $\mathcal S \subseteq \mathcal T$
implies $G(\mathcal S) \le G(\mathcal T)$. We allocate $G(\mathcal G)$ by the Shapley value
\citep{shapley1953value},
\begin{equation}
  \phi_g = \sum_{\mathcal S \subseteq \mathcal G \setminus \{g\}}
  \frac{|\mathcal S|!\,(|\mathcal G|-|\mathcal S|-1)!}{|\mathcal G|!}
  \big[G(\mathcal S \cup \{g\}) - G(\mathcal S)\big],
  \qquad s_g = \frac{\phi_g}{G(\mathcal G)},
  \label{eq:shapley}
\end{equation}
the average marginal gain of group $g$ over all orders of entry. It is the only allocation that is efficient
($\sum_g \phi_g = G(\mathcal G)$), symmetric, assigns zero to a group that never adds fit, and is additive in $G$; by
monotonicity, every $s_g \in [0,1]$. Shapley allocations of fit are standard for correlated regressors
\citep{lipovetsky2001analysis,huettner2012axiomatic}. To report coarser groups, such as \emph{fit} = quality $\cup$
demand match, we apply \eqref{eq:shapley} to the game on blocks $v(\mathcal T) = G(\cup \mathcal T)$, whose values
are among the \nShapFits{} fits already computed. Shares are reported in-sample and on the holdout; on the holdout a
group gains only by predicting future choices, and $G_h$ need not be monotone.

% ---------------------------------------------------------------------------------------------------------------------
\subsection{Bounding general quality}
\label{sec:bounds}

Benchmark scores cover few models, so a small quality share could reflect missing measurement rather than
indifference to quality. We bound it from above with model-specific intercepts for the most used models.

\begin{definition}[Hubs]
\label{def:hubs}
Let $t_h$ be the train--holdout cut (\tHold). The \emph{hub set} $\mathcal H$ is the 500 models with the highest
degree $d_l(t_h)$ over the whole Hub (all modalities; ties broken by model index). $\mathcal H$ uses only training
data and does not depend on the sampling design.
\end{definition}

The fixed-effect model adds an intercept for each hub,
\begin{equation}
  V^{\mathrm{FE}}_{kl} = \theta^\top x_{kl} + \sum_{h \in \mathcal H} \phi_h \1[l = h],
  \label{eq:fe}
\end{equation}
estimated by \eqref{eq:corrected} with the same offset (\cref{prop:pps} holds for any $V$) and a small ridge penalty
$\frac{\lambda}{2}\|\phi\|^2$, $\lambda = 10^{-4}$. Non-hubs have $\phi = 0$, so $\phi_h$ is hub $h$'s
time-invariant attractiveness relative to the rest of the market. It absorbs every fixed attribute of the hub:
unmeasured quality, but also brand, tutorials and default status. Entering the fixed effects as an additional player
in \cref{def:game} gives a share $s_{\mathrm{FE}}$, and
\[
  \bar s_{\mathrm{qual}} = s_{\mathrm{qual}} + s_{\mathrm{FE}}
\]
is the largest share that any time-invariant hub attribute could contribute to quality (\feShare\% of the gain
in-sample).

To see how much of the fixed effects is quality, we project the estimated intercepts on external quality measures
$z_h$ (Open LLM Leaderboard v2 average; arena ratings) with controls $w_h$ (parameter count, re-distributor accounts),
$\hat\phi_h = b_0 + b^\top z_h + \gamma^\top w_h + u_h$, over the \extHubs{} hubs that are covered and precisely
estimated.

\begin{lemma}[Noise-corrected quality share]
\label{lem:noise}
Suppose $\hat\phi_h = \phi_h + e_h$ with $\E[e_h \mid \phi_h, z_h, w_h] = 0$ and $\mathrm{Var}(e_h) = \sigma_h^2$.
Then $\mathrm{Var}(\hat\phi) = \mathrm{Var}(\phi) + \bar\sigma^2$, the projection coefficients are unaffected by $e$,
and
\[
  q = \frac{\mathrm{Var}(\hat b^\top z_h)}{\mathrm{Var}(\hat\phi_h) - \bar\sigma^2}
\]
estimates the share of $\mathrm{Var}(\phi)$ explained by external quality, with $\sigma_h^2$ the bootstrap variance
of $\hat\phi_h$.
\end{lemma}

\noindent We split the fixed-effect share proportionally, $s^{\mathrm{ext}}_{\mathrm{qual}} = s_{\mathrm{qual}} + q\,
s_{\mathrm{FE}}$ ($q = \extQ\%$). Shapley shares are not linear in a split of a player, so this split is an
approximation; the two end points, $s_{\mathrm{qual}}$ and $\bar s_{\mathrm{qual}}$, do not depend on it
(\cref{app:proofs}).

% ---------------------------------------------------------------------------------------------------------------------
\subsection{Does the decomposition recover known mechanisms?}
\label{sec:validation}

The decomposition attributes most of the likelihood gain to popularity (\cref{sec:results}). Because the groups are
correlated, this could in principle be a property of the method rather than of the data. We test the method on
choices whose generating mechanism is known by construction.

For a generating utility $V^{(r)}$, keep every real candidate set $\D_k$ with its real features and inclusion
probabilities, and replace the observed parent $j_k$ by a draw
\begin{equation}
  \tilde j^{(r)}_k \sim P^{(r)}(l \mid \D_k) = \frac{\exp(V^{(r)}_{kl} - \log \pi_{kl})}{\sum_{m \in \D_k}
  \exp(V^{(r)}_{km} - \log \pi_{km})}, \qquad l \in \D_k,
  \label{eq:sir}
\end{equation}
the model-implied choice on the sampled set (\cref{prop:pps}). We use three utilities built from the estimates of
\cref{sec:estimation,sec:bounds}:
\[
  V^{\mathrm{SF}}_{kl} = \hat\theta_{\mathrm{pop}}^\top x_{kl,\mathrm{pop}}, \qquad
  V^{\mathrm{QF}}_{kl} = \hat\theta_{\mathrm{qual}}^\top x_{kl,\mathrm{qual}} + \hat\theta_{\mathrm{match}}^\top
  x_{kl,\mathrm{match}} + \hat q_l, \qquad
  V^{\mathrm{FULL}}_{kl} = \hat V^{\mathrm{FE}}_{kl},
\]
where $\hat q_l = \hat b^\top z_l$ is the part of the hub intercept explained by external quality ($\hat q_l = 0$ for
non-hubs). SF lets choices depend only on popularity, QF only on quality and demand match, and FULL on the estimated
model. Running steps 2--3 of Algorithm 1 on $\{(\tilde j^{(r)}_k, \D_k)\}_k$ gives shares $\hat s(r)$. If the
decomposition separates the mechanisms, $\hat s_{\mathrm{pop}}(\mathrm{SF})$ and $\hat s_{\mathrm{fit}}(\mathrm{QF})$
should be close to one, and $\hat s(\mathrm{FULL})$ should equal the shares estimated on the real choices; $1 -
\hat s_{\mathrm{pop}}(\mathrm{SF})$ and $1 - \hat s_{\mathrm{fit}}(\mathrm{QF})$ measure how much fit leaks to a
correlated group.

The decomposition attributes \recSfPop\% of the gain to popularity under SF and \recQfFit\% to fit under QF, and it
returns \recFullPop\%, \recFullMatch\% and \recFullInertia\% (popularity, demand match, self-supply plus habit) under
FULL, the shares estimated on the real choices.\todo{Recovery is run with four groups; rerun with five or keep the
note that self-supply + habit = the former inertia group.} When choices are generated by quality and demand match
alone, the decomposition does not manufacture a popularity share. As a check of the whole pipeline, we also grow
synthetic networks in which a new node attaches to $l$ with weight $d_l^{\alpha}$ if $d_l > 0$ and $e^{\beta_0}$
otherwise ($\alpha = 1$, $\beta_0 = -1$), replay them with the same code, and recover $\hat\alpha$ within three
standard errors of one under all three sampling designs (\cref{app:validation}).

% ---------------------------------------------------------------------------------------------------------------------
\subsection{Samples, weights and author types}
\label{sec:samples}

The estimator above is defined for any set of occasions and weights. The main specification uses first adoptions
with $w_k = 1$; \cref{sec:results} reports the alternatives defined here, and a rule for calling the results robust
that we fixed before estimating them.

\paragraph{Samples.} \emph{P}: first adoptions (main). \emph{D}: all 24-hour decisions, including repeat decisions.
\emph{U}: all uploads. In D and U the author's earlier parents enter $\D_k$ with $\pi = 1$, so \cref{prop:pps}
applies unchanged. On a repeat decision the degree of a candidate includes the author's own earlier children of it,
which would load habit onto popularity. We therefore replace $d_{kl}$ by $\tilde d_{kl} = d_{kl} - o_{kl}$, where
$o_{kl}$ counts the children of $l$ by author $a_k$ created before $t_k$ (and similarly for adoptions in the last 30
days). On P, $o_{kl} \equiv 0$.

\paragraph{Weights.} With $N$ occasions by $A$ authors and $n_a$ occasions of author $a$, \emph{author weights}
$w_k = (N/A)/n_{a_k}$ give every author the same total weight and a mean weight of one. They estimate the behavior of
the typical author rather than of the typical choice.

\paragraph{Author types.} Authors are classified into six types (major organizations, re-distribution accounts,
other organizations, and individuals with $\ge$100, 10--99 and $<$10 models; \cref{app:data}). Heterogeneity is
tested by $LR = 2\big(\sum_\tau \ell_\tau(\hat\theta_\tau) - \ell(\hat\theta)\big)$, asymptotically $\chi^2$ with
degrees of freedom equal to the number of additional parameters of the type-specific models; because the logit scale may differ across types, types are compared by
shares, not coefficients.

\paragraph{Robustness rule.} An alternative sample or weighting is called robust if its popularity share is within 10
percentage points of the main specification and its fit/popularity ratio within a factor of three.
, and add
% % Proofs for sections/4_method_revised.tex. Input after the other appendices when the revised Section 4 is used.
\section{Proofs for Section~\ref{sec:method}}
\label{app:proofs}

\paragraph{Proof of \cref{lem:size}.}
(i) Let $I_{kl} = \1[l \in \D_k]$. For $l \ne j_k$, $I_{kl} \sim \mathrm{Bern}(\pi_{kl})$; $I_{kj_k} = 1$ because
the chosen parent is added if it was not drawn. By linearity of expectation (independence is not needed),
$\E|\D_k| = \sum_{l \ne j_k} \pi_{kl} + 1 = \sum_{l} \pi_{kl} + (1 - \pi_{kj_k})$, which lies in $[S, S+1]$ when
$\sum_l \pi_{kl} = S$.

(ii) Each term $n_\delta \min\{1, c(\delta+1)\}$ is continuous and non-decreasing in $c$, linear with slope
$n_\delta(\delta+1) > 0$ for $c < 1/(\delta+1)$ and constant afterwards; hence $\Sigma$ is continuous,
non-decreasing, strictly increasing wherever at least one class is unsaturated, $\Sigma(0) = 0$ and $\Sigma(c) =
|\M_k|$ once every class is saturated (for $c \ge 1$). If $|\M_k| > S$, the intermediate-value theorem gives a
solution, and at any solution some class is unsaturated (otherwise $\Sigma = |\M_k| > S$), so $\Sigma$ is strictly
increasing there and the solution is unique. Since $1/(\delta+1)$ decreases in $\delta$, the saturated classes at
$c_k$ are always the $m$ largest degrees. Given $m$, $\Sigma(c) = \sum_{i \le m} n_{\delta_i} + c\,(n_0 + \sum_{i >
m} n_{\delta_i}(\delta_i + 1))$, and solving $\Sigma(c) = S$ gives the formula. The consistency conditions select
the true $m$; trying $m = 0, 1, \dots$ in turn finds it in one pass over the distinct degrees, and the replay keeps
the degree histogram up to date incrementally. Degree-0 models are drawn as a $\mathrm{Bin}(n_0, c_k)$ count followed
by rejection sampling. \hfill$\square$

\paragraph{Proof of \cref{prop:pps}.}
Fix $k$ and drop the index. Every $l \ne j$ enters $\D$ independently with probability $\pi_l$, and $j$ is always in
$\D$, so given that $j$ was chosen the probability of a set $\D \ni j$ is $q(\D \mid j) = \prod_{l \in \D \setminus
\{j\}} \pi_l \prod_{l \notin \D}(1 - \pi_l)$. For $j, l \in \D$ the two products differ only in the factors of $j$
and $l$, so $q(\D \mid l) = q(\D \mid j)\,\pi_j / \pi_l$. By Bayes' rule and \eqref{eq:logit},
\[
  P(j \mid \D) = \frac{P(j \mid \M)\, q(\D \mid j)}{\sum_{l \in \D} P(l \mid \M)\, q(\D \mid l)}
  = \frac{e^{V_j}\, q(\D \mid j)}{\sum_{l \in \D} e^{V_l}\, q(\D \mid j)\, \pi_j / \pi_l}
  = \frac{e^{V_j - \log \pi_j}}{\sum_{l \in \D} e^{V_l - \log \pi_l}}.
\]
Only $P(j \mid \M) \propto e^{V_j}$ is used, so the result holds for any $V$ (in particular \eqref{eq:fe}).
\hfill$\square$

\paragraph{Proof of \cref{cor:consistency} (sketch).}
By \cref{prop:pps}, $\ell_k(\theta) = \log P(j_k \mid \D_k; \theta)$ is the log of a correctly specified conditional
density of the observed $j_k$ given $\D_k$, so $\theta_0$ maximizes $\E[\ell_k(\theta)]$ (information inequality).
By \cref{lem:concave}, $\ell_k$ is concave, and $-\nabla^2 \E[\ell_k] = \E[\mathrm{Cov}_{p_k}(x_k)]$ is positive
definite at $\theta_0$, so the maximizer is unique. Group occasions by author: the author sums
$L_a(\theta) = \sum_{k: a_k = a} \ell_k(\theta)$ are independent with bounded size, so a law of large numbers holds
pointwise, and for concave objectives pointwise convergence to a uniquely maximized limit implies consistency
\citep[Thm.~2.7]{newey1994large}. Asymptotic normality follows from a mean-value expansion of the first-order
condition and the central limit theorem for the independent author scores $\nabla L_a(\theta_0)$, giving the
sandwich $H_0^{-1} \Omega_0 H_0^{-1}$ \citep[Thm.~3.1]{newey1994large}. \hfill$\square$

\paragraph{Proof of \cref{lem:concave}.}
For one occasion, $\ell_k = v_{kj_k} - \log \sum_l e^{v_{kl}}$ with $\partial v_{kl} / \partial \theta = x_{kl}$.
Then $\nabla \ell_k = x_{kj_k} - \sum_l p_{kl} x_{kl}$, and differentiating $\bar x_k = \sum_l p_{kl} x_{kl}$ with
$\partial p_{kl} / \partial \theta = p_{kl}(x_{kl} - \bar x_k)$ gives $\nabla^2 \ell_k = -\sum_l p_{kl}(x_{kl} -
\bar x_k)(x_{kl} - \bar x_k)^\top = -\mathrm{Cov}_{p_k}(x_k)$, a negative semi-definite matrix. Summing with
non-negative weights preserves this; if the weighted sum of covariances is positive definite at every $\theta$, the
Hessian is negative definite and the maximizer is unique. \hfill$\square$

\paragraph{Proof of \cref{lem:noise}.}
With $\hat\phi_h = \phi_h + e_h$ and $\E[e_h \mid \phi_h, z_h, w_h] = 0$, $\mathrm{Cov}(\phi_h, e_h) = 0$, so
$\mathrm{Var}(\hat\phi) = \mathrm{Var}(\phi) + \E[\sigma_h^2]$. The population projection of $\hat\phi$ on $(z, w)$
equals that of $\phi$, because $\mathrm{Cov}(z, e) = \mathrm{Cov}(w, e) = 0$; hence $\mathrm{Var}(b^\top z)$ is the
variance of the quality part of $\phi$ itself, and dividing it by $\mathrm{Var}(\hat\phi) - \bar\sigma^2$ estimates
its share of $\mathrm{Var}(\phi)$. Without the correction the denominator is too large and the quality share too
small. We use the bootstrap variance of each $\hat\phi_h$ for $\sigma_h^2$ and keep hubs with bootstrap standard error
below one. \hfill$\square$

\paragraph{The proportional split.}
$s^{\mathrm{ext}}_{\mathrm{qual}} = s_{\mathrm{qual}} + q\, s_{\mathrm{FE}}$ assumes that the quality and the
non-quality parts of the hub intercepts contribute to the likelihood gain in proportion to their variances. An exact
allocation would enter the quality part $\hat b^\top z_h$ as a separate Shapley player, which requires external
quality for every hub; coverage is limited to \extHubs{} hubs. The lower end ($s_{\mathrm{qual}}$, no fixed effects)
and the upper end ($\bar s_{\mathrm{qual}}$, the whole fixed-effect share counted as quality) do not use the split.

\paragraph{The recovery draws.}
In \eqref{eq:sir}, $\D_k$ is the realized candidate set, which was drawn with the real $j_k$ forced in. The
simulated choice is therefore a draw from the model-implied conditional law on the realized set; the forced
inclusion of one real parent among about 300 candidates has a negligible effect on the composition of the set.
The draws use the four-group decomposition of the earlier analysis, whose inertia group equals self-supply and habit
combined; the full model, and hence $V^{\mathrm{FULL}}$, is the same under four and five groups.
 next to the other appendices (it holds the proofs cited here).
% All labels used elsewhere in the paper (sec:method, eq:logit, sec:sampling, eq:pps, prop:pps, eq:corrected,
% sec:shapley, eq:shapley, sec:bounds, sec:validation) are kept.
%
% Theorem environments: acmart defines definition, lemma, proposition, corollary; assumption and remark are
% defined below if missing (no change to main.tex needed).
%
% Changes relative to 4_method.tex, by comment of the author:
%  1. 4.1 now defines the choice occasion, degree and the information set (Def. 1) before the random-utility
%     assumption (Assumption 1); IIA is stated, because Prop. 1 uses it.
%  2. Consistency / asymptotic normality is a separate corollary (Cor. 1) after the sampling proposition.
%  3. 4.3 is about the estimator: weighted sampled likelihood, gradient and Hessian, concavity (Lemma 3), the
%     estimation algorithm (Algorithm 1), the clustered sandwich and the multiplier bootstrap. Hardware is one
%     sentence.
%  4. 4.4 starts from why coefficients cannot answer the question (units, groups, within-set variation, order
%     dependence with the real numbers) and defines the likelihood-gain game (Def. 3) before the Shapley value.
%  5. 4.5 defines the hub set formally (Def. 4) and states the noise-corrected split as a lemma.
%  6. 4.6 writes the three generating utilities, the simulated choice and the recovery criterion as equations.
%  New 4.7: samples, weights and author types (definitions only; results in Section 5, Table tab:samples).
%  Five groups (self-supply and habit replace inertia) and 2^5 = 32 fits throughout.
%
% Section 3 needs two matching edits (not done): (a) "the 500 most adopted parents" -> refer to Def. 4 (the hubs are
% ranked by degree d_l(t_h) over all children, not by first adoptions); (b) the inertia group -> self-supply
% (same_author) and habit (author_reuse, author_family, is_last).
% =====================================================================================================================
\makeatletter
\@ifundefined{assumption}{\newtheorem{assumption}{Assumption}}{}
\@ifundefined{remark}{\theoremstyle{acmdefinition}\newtheorem{remark}[theorem]{Remark}\theoremstyle{acmplain}}{}
\makeatother

\section{Model and Estimation}
\label{sec:method}

\Cref{sec:data} showed that adoption is highly concentrated and that popularity, general quality and demand match
could each produce this concentration. To separate them we need a model in which all three can act at once and be
compared on a common scale. This section builds that model (\cref{sec:choice}), makes it computable on markets of
hundreds of thousands of alternatives (\cref{sec:sampling,sec:estimation}), turns its estimates into a
decomposition by signal group (\cref{sec:shapley}), bounds what unmeasured quality could add (\cref{sec:bounds}),
and checks that the decomposition recovers known mechanisms (\cref{sec:validation}).

% ---------------------------------------------------------------------------------------------------------------------
\subsection{Base-model selection as a discrete choice}
\label{sec:choice}

\begin{definition}[Choice occasion]
\label{def:occasion}
A \emph{choice occasion} $k$ is a first adoption (\cref{sec:data}), described by its author $a_k$, time $t_k$,
modality $m_k$ and chosen parent $j_k$. Its \emph{market} $\M_k$ is the set of models of modality $m_k$ created
before $t_k$ that $a_k$ has not used as a parent. The \emph{degree} of a model $l$ at time $t$ is
\[
  d_l(t) = \#\{c : \text{$c$ declares $l$ as its parent with type fine-tune, adapter or unspecified, and $c$ was
  created before $t$}\},
\]
and $d_{kl} = d_l(t_k)$. Every candidate $l \in \M_k$ carries a feature vector $x_{kl} \in \mathbb{R}^p$ computed
only from models, edges and metadata that exist before $t_k$.
\end{definition}

The features are the five groups $g \in \mathcal G = \{\text{pop}, \text{qual}, \text{match}, \text{self},
\text{habit}\}$ of \cref{sec:data}: popularity, general quality, demand match, self-supply (the candidate is the
author's own model) and habit (the candidate belongs to a family the author used before). We write $x_{kl} =
(x_{kl,g})_{g \in \mathcal G}$ and $\theta = (\theta_g)_{g \in \mathcal G}$ conformably. Popularity and quality
features vary over candidates and time; demand-match, self-supply and habit features also vary over authors.

\begin{assumption}[Random utility]
\label{ass:rum}
The author of occasion $k$ chooses $j_k = \arg\max_{l \in \M_k} U_{kl}$, where
\[
  U_{kl} = V_{kl} + \varepsilon_{kl}, \qquad V_{kl} = \theta^\top x_{kl},
\]
and the $\varepsilon_{kl}$ are i.i.d.\ standard Type-I extreme value, independent of $\{x_{kl}\}_{l \in \M_k}$.
\end{assumption}

\begin{lemma}[\citealp{mcfadden1974conditional}]
\label{lem:logit}
Under \cref{ass:rum},
\begin{equation}
  P(j \mid \M_k) = \frac{\exp(\theta^\top x_{kj})}{\sum_{l \in \M_k} \exp(\theta^\top x_{kl})}.
  \label{eq:logit}
\end{equation}
\end{lemma}

\begin{remark}[What the model does and does not say]
\label{rem:logit}
(i) \emph{Identification.} Only differences across candidates within an occasion enter \eqref{eq:logit}; a
variable that is constant across $l \in \M_k$ (author experience, organizational status) cancels. Author
heterogeneity therefore enters through author-specific features (demand match, self-supply, habit) and through
estimation by author type (\cref{sec:samples}).
(ii) \emph{Scale.} The error variance is normalized, so $\theta$ is measured relative to the dispersion of tastes;
comparisons across samples are made with likelihood shares (\cref{sec:shapley}), which are scale free, not with
coefficients.
(iii) \emph{Independence of irrelevant alternatives.} $P(j \mid \M_k) / P(l \mid \M_k) = \exp(V_{kj} - V_{kl})$
does not depend on the other alternatives. \Cref{prop:pps} uses this property.
(iv) \emph{Interpretation.} We read $\theta$ descriptively, as the weights that best summarize which signals
choices co-move with, not as causal effects (\cref{sec:discussion}).
\end{remark}

Two network-growth models are special cases. With $\theta = 0$, \eqref{eq:logit} is \emph{uniform attachment}:
every model in the market is equally likely. With only $x_{kl} = (\log \max(d_{kl},1),\ \1[d_{kl}>0])$ and
coefficients $(\alpha, \beta)$,
\[
  P(j \mid \M_k) \;\propto\;
  \begin{cases} e^{\beta}\, d_{kj}^{\alpha}, & d_{kj} > 0, \\ 1, & d_{kj} = 0, \end{cases}
\]
which is \emph{preferential attachment} with exponent $\alpha$; $\alpha = 1$ is the linear model of
\citet{barabasi1999emergence}. Following \citet{overgoor2019choosing}, \eqref{eq:logit} nests these mechanisms and
lets them be compared by likelihood together with quality and demand match.

% ---------------------------------------------------------------------------------------------------------------------
\subsection{Sampling the competitors that matter}
\label{sec:sampling}

The denominator of \eqref{eq:logit} has a median of \marketMedian{} terms, and the training data contain \nTrain{}
occasions, so the likelihood cannot be evaluated on full markets. We evaluate it on a sampled candidate set.

\begin{definition}[Degree-proportional Poisson design]
\label{def:pps}
For occasion $k$ and a target size $S$, every $l \in \M_k$ enters the candidate set $\D_k$ independently with
probability
\begin{equation}
  \pi_{kl} = \min\{1,\; c_k (d_{kl} + 1)\},
  \label{eq:pps}
\end{equation}
where the scale $c_k > 0$ solves $\sum_{l \in \M_k} \pi_{kl} = S$. The chosen parent $j_k$ is added to $\D_k$ if it
was not drawn.
\end{definition}

The inclusion probability is proportional to $d+1$ (the $+1$ keeps never-used models in the design), and $c_k$ is
the factor of proportionality. The condition on $c_k$ fixes the expected size of the set:

\begin{lemma}[Set size and the scale $c_k$]
\label{lem:size}
(i) $\E|\D_k| = \sum_{l \in \M_k} \pi_{kl} + (1 - \pi_{kj_k}) \in [S, S+1]$.
(ii) Let $n_\delta$ be the number of models in $\M_k$ with degree $\delta$, and let $\delta_1 > \delta_2 > \dots >
\delta_M > 0$ be the positive degrees. The function $\Sigma(c) = \sum_\delta n_\delta \min\{1, c(\delta+1)\}$ is
continuous and non-decreasing, strictly increasing wherever some degree class is unsaturated
($c(\delta+1) < 1$), and runs from $0$ to $|\M_k|$. If $|\M_k| > S$, $\Sigma(c) = S$ has a unique solution
\[
  c_k = \frac{S - \sum_{i \le m} n_{\delta_i}}{n_0 + \sum_{i > m} n_{\delta_i}(\delta_i + 1)},
\]
where $m \in \{0, \dots, M\}$ is the number of saturated classes: the unique $m$ with $c_k(\delta_i + 1) \ge 1$ for
$i \le m$, $c_k(\delta_i + 1) < 1$ for $i > m$, and $c_k \le 1$.
\end{lemma}

\noindent The proof is in \cref{app:proofs}. With $S = 300$, on average \piOneMean{} candidates per training
occasion have $\pi = 1$; these are the incumbents a developer is most likely to weigh, and \piOneChosen\% of chosen
parents are among them. The rest of each set is a probability-proportional-to-size sample of the long tail.

\begin{proposition}[Sampling correction]
\label{prop:pps}
Under \cref{ass:rum} and \cref{def:pps},
\begin{equation}
  P(j \mid \D_k) = \frac{\exp(\theta^\top x_{kj} - \log \pi_{kj})}{\sum_{l \in \D_k} \exp(\theta^\top x_{kl} - \log
  \pi_{kl})}, \qquad j \in \D_k.
  \label{eq:corrected}
\end{equation}
\end{proposition}

\noindent Given that $j$ was chosen, the probability of drawing a set $\D \ni j$ is $q(\D \mid j) = \prod_{l \in \D
\setminus \{j\}} \pi_l \prod_{l \notin \D} (1 - \pi_l)$, so $q(\D \mid j)/q(\D \mid l) = \pi_l / \pi_j$ for any $j,
l \in \D$; Bayes' rule and \eqref{eq:logit} give \eqref{eq:corrected} (\cref{app:proofs};
\citealp{mcfadden1978modelling}). The argument uses only the logit form of $P(j \mid \M_k)$, so it holds for any
specification of $V_{kl}$, including the fixed-effect model of \cref{sec:bounds}.

\begin{corollary}[Consistency and asymptotic normality]
\label{cor:consistency}
Let $\hat\theta$ maximize $\ell(\theta) = \sum_k \log P(j_k \mid \D_k; \theta)$ from \eqref{eq:corrected}. Suppose
occasions are independent across authors with bounded numbers per author, the parameter space is compact with the
true $\theta_0$ in its interior, $\E\|x_{kl}\|^2 < \infty$, and $H_0 = \E[\mathrm{Cov}_{p_k}(x_k)]$ is positive
definite, where $\mathrm{Cov}_{p_k}$ is the covariance of $x_{kl}$ over $l \in \D_k$ under the probabilities
\eqref{eq:corrected}. Then, as the number of authors $A \to \infty$, $\hat\theta \to \theta_0$ in probability and
$\sqrt{A}(\hat\theta - \theta_0) \to N(0, H_0^{-1} \Omega_0 H_0^{-1})$, with $\Omega_0$ the variance of an author's
summed score.
\end{corollary}

\noindent By \cref{prop:pps}, $\ell$ is a correctly specified conditional likelihood of the observed $(j_k, \D_k)$,
and it is concave (\cref{lem:concave} below); the result follows from the standard argument for concave M-estimators
\citep{newey1994large}. Two consequences of \eqref{eq:corrected} matter for what follows.

\begin{remark}[Uniform attachment as the baseline, and why sample by degree]
\label{rem:pps}
(i) At $\theta = 0$, \eqref{eq:corrected} gives each sampled candidate the weight $1/\pi_{kl}$, and $\sum_{l \in
\D_k} 1/\pi_{kl}$ is a Horvitz--Thompson estimate of $|\M_k|$ \citep{horvitz1952generalization}; uniform attachment
on the sample therefore corresponds to uniform attachment on the full market. It is the baseline of every likelihood
gain below.
(ii) \Cref{cor:consistency} holds for any design with $\pi > 0$, but the information matrix $H_0 = \E[\mathrm{Cov}_{p_k}
(x_k)]$ depends on the design: a candidate set contributes information only through the variation of $x$ among the
candidates that have non-negligible choice probability. A uniform sample of 300 from $3 \times 10^5$ models almost
never contains the few models that receive most adoptions, so the chosen parent is compared mostly with never-used
models. Under a uniform design the median clustered standard error is \seRatioUnif{} times larger, and the
estimated weights themselves change (\cref{sec:robustq}).
\end{remark}

% ---------------------------------------------------------------------------------------------------------------------
\subsection{Estimation and inference}
\label{sec:estimation}

\Cref{cor:consistency} defines the estimator; we now state how it is computed and how its uncertainty is measured.
Write $v_{kl}(\theta) = \theta^\top x_{kl} - \log \pi_{kl}$, $p_{kl}(\theta) = e^{v_{kl}} / \sum_{m \in \D_k}
e^{v_{km}}$ and $\bar x_k(\theta) = \sum_{l \in \D_k} p_{kl} x_{kl}$. With occasion weights $w_k$ ($w_k = 1$ in the
main specification; other weights in \cref{sec:samples} and the bootstrap below), the sampled log-likelihood is
\begin{equation}
  \ell(\theta) = \sum_k w_k \Big[ v_{kj_k}(\theta) - \log \sum_{l \in \D_k} e^{v_{kl}(\theta)} \Big].
  \label{eq:loglik}
\end{equation}

\begin{lemma}[Concavity]
\label{lem:concave}
$\nabla \ell(\theta) = \sum_k w_k (x_{kj_k} - \bar x_k)$ and $\nabla^2 \ell(\theta) = -\sum_k w_k
\mathrm{Cov}_{p_k}(x_k)$. Hence $\ell$ is concave, and strictly concave, with a unique maximizer, whenever
$\sum_k w_k \mathrm{Cov}_{p_k}(x_k)$ is positive definite.
\end{lemma}

\noindent The offset $-\log \pi$ does not depend on $\theta$, so it changes neither the gradient's form nor the
concavity. The first-order condition equates the weighted mean features of the chosen parents with their mean under
the model. Concavity has a practical consequence: every refit in the decomposition and the bootstrap converges to
the same optimum from any starting point, so warm starts change only computing time.

\begin{center}
\fbox{\begin{minipage}{0.94\linewidth}
\textbf{Algorithm 1 (estimation).}
\begin{enumerate}
\item \emph{Replay.} Process the network in time order; at each occasion $k$ update $d_l(t_k)$, solve $c_k$
  (\cref{lem:size}), draw $\D_k$ by \cref{def:pps}, and record $(x_{kl}, \log \pi_{kl})_{l \in \D_k}$.
\item \emph{Fit.} Maximize \eqref{eq:loglik} by L-BFGS with the analytic gradient of \cref{lem:concave}; compute
  the Hessian at $\hat\theta$.
\item \emph{Subsets.} For every $\mathcal S \subseteq \mathcal G$, maximize \eqref{eq:loglik} with $\theta_g = 0$ for
  $g \notin \mathcal S$, warm-started from the fit of $\mathcal S$ minus one group ($2^{|\mathcal G|} = \nShapFits{}$
  fits; \cref{sec:shapley}).
\item \emph{Bootstrap.} Repeat steps 2--3 $B = 50$ times with weights $w_k \xi_{a_k}$, $\xi_a \sim \mathrm{Exp}(1)$
  i.i.d.\ over authors, warm-started from the point estimates.
\end{enumerate}
\end{minipage}}
\end{center}

\noindent The training data have \nRowsTrain{} million (occasion, candidate) rows; the full model fits in about a
minute on one GPU.\footnote{Candidate sets are grouped into buckets of similar size and kept on the GPU in single
precision; the likelihood, gradient and Hessian are accumulated in double precision.}

\paragraph{Inference.}
Occasions of one author are dependent (the same project, the same tastes, batch uploads): the \nFirst{} first
adoptions come from \nFirstAuthors{} authors, and the median clustered standard error is \seClusterInfl{} times
the one that treats occasions as independent. We therefore use the cluster-robust sandwich
\[
  \widehat{\mathrm{Var}}(\hat\theta) = \hat H^{-1} \Big( \textstyle\sum_a \hat s_a \hat s_a^\top \Big) \hat H^{-1},
  \qquad \hat H = -\nabla^2 \ell(\hat\theta), \quad \hat s_a = \sum_{k: a_k = a} w_k (x_{kj_k} - \bar x_k(\hat\theta)),
\]
\citep{cameron2015practitioner}. For derived quantities, the Shapley shares of \cref{sec:shapley} in particular, we
use the multiplier bootstrap of step 4: replicate $b$ maximizes $\sum_k \xi^{(b)}_{a_k} w_k \ell_k(\theta)$; the
shares are recomputed in each replicate; and intervals are the 2.5\% and 97.5\% percentiles. Normalized Exp(1)
weights are Dirichlet(1,\dots,1), the Bayesian bootstrap of \citet{rubin1981bayesian}; reweighting whole authors
preserves within-author dependence, and the weighted bootstrap is valid for M-estimators of this kind
\citep{ma2005robust}.

\paragraph{Out-of-sample evaluation.}
Let $K_h$ be the holdout occasions (the last 10\% in time) and $P_0$ uniform attachment ($\theta = 0$ in
\eqref{eq:corrected}). We report the mean holdout log-likelihood gain
\[
  \Delta_h = \frac{1}{|K_h|} \sum_{k \in K_h} \big[ \log P(j_k \mid \D_k; \hat\theta) - \log P_0(j_k \mid \D_k)
  \big],
\]
so that $e^{\Delta_h}$ is the geometric-mean factor by which the model raises the probability of the true parent
over uniform attachment ($\Delta_h = \holdGain$, a factor of \holdGainTimes), and the top-$k$ rate, the share of
holdout occasions in which $j_k$ ranks among the $k$ candidates with the highest $\hat V_{kl}$ (ranking by
utility, without the offset; ties broken at random).

% ---------------------------------------------------------------------------------------------------------------------
\subsection{From coefficients to a decomposition}
\label{sec:shapley}

$\hat\theta$ has 26 components, but the question is about groups of signals: how much of developers' choices do
popularity signals account for, and how much do quality and demand match? The coefficients cannot answer it
directly, for four reasons.
(i) Features have different units (indicators, logarithms of counts, similarities), so coefficients are not
comparable across features.
(ii) Each group contains several features, and the question concerns the group.
(iii) What a group contributes to explaining a choice depends on how much its \emph{index}
$I_{g,kl} = \theta_g^\top x_{kl,g}$ varies across the candidates of $\D_k$, not on the size of $\theta_g$ alone;
a large coefficient on a feature that rarely varies within a set explains little.
(iv) The groups are correlated: popular models are also well documented and match common tasks. The fit gained by
adding a group therefore depends on which groups are already in the model. Demand match accounts for \matchFirst\%
of the full model's likelihood gain when it enters first and for \matchLast\% when it enters last.
We therefore measure each group by its contribution to the model's fit and allocate that contribution with a rule
that does not depend on the order of entry.

\begin{definition}[Likelihood-gain game]
\label{def:game}
For $\mathcal S \subseteq \mathcal G$, let $\hat\theta_{\mathcal S}$ maximize \eqref{eq:loglik} subject to
$\theta_g = 0$ for $g \notin \mathcal S$, and let
\[
  G(\mathcal S) = \ell(\hat\theta_{\mathcal S}) - \ell(0),
\]
the log-likelihood gain over uniform attachment of the model with the groups in $\mathcal S$. The pair
$(\mathcal G, G)$ is a cooperative game with $G(\emptyset) = 0$. Its holdout version $G_h(\mathcal S)$ evaluates the
same $\hat\theta_{\mathcal S}$ by the mean holdout log-likelihood gain.
\end{definition}

Because the models are nested and fitted by maximum likelihood, $G$ is monotone: $\mathcal S \subseteq \mathcal T$
implies $G(\mathcal S) \le G(\mathcal T)$. We allocate $G(\mathcal G)$ by the Shapley value
\citep{shapley1953value},
\begin{equation}
  \phi_g = \sum_{\mathcal S \subseteq \mathcal G \setminus \{g\}}
  \frac{|\mathcal S|!\,(|\mathcal G|-|\mathcal S|-1)!}{|\mathcal G|!}
  \big[G(\mathcal S \cup \{g\}) - G(\mathcal S)\big],
  \qquad s_g = \frac{\phi_g}{G(\mathcal G)},
  \label{eq:shapley}
\end{equation}
the average marginal gain of group $g$ over all orders of entry. It is the only allocation that is efficient
($\sum_g \phi_g = G(\mathcal G)$), symmetric, assigns zero to a group that never adds fit, and is additive in $G$; by
monotonicity, every $s_g \in [0,1]$. Shapley allocations of fit are standard for correlated regressors
\citep{lipovetsky2001analysis,huettner2012axiomatic}. To report coarser groups, such as \emph{fit} = quality $\cup$
demand match, we apply \eqref{eq:shapley} to the game on blocks $v(\mathcal T) = G(\cup \mathcal T)$, whose values
are among the \nShapFits{} fits already computed. Shares are reported in-sample and on the holdout; on the holdout a
group gains only by predicting future choices, and $G_h$ need not be monotone.

% ---------------------------------------------------------------------------------------------------------------------
\subsection{Bounding general quality}
\label{sec:bounds}

Benchmark scores cover few models, so a small quality share could reflect missing measurement rather than
indifference to quality. We bound it from above with model-specific intercepts for the most used models.

\begin{definition}[Hubs]
\label{def:hubs}
Let $t_h$ be the train--holdout cut (\tHold). The \emph{hub set} $\mathcal H$ is the 500 models with the highest
degree $d_l(t_h)$ over the whole Hub (all modalities; ties broken by model index). $\mathcal H$ uses only training
data and does not depend on the sampling design.
\end{definition}

The fixed-effect model adds an intercept for each hub,
\begin{equation}
  V^{\mathrm{FE}}_{kl} = \theta^\top x_{kl} + \sum_{h \in \mathcal H} \phi_h \1[l = h],
  \label{eq:fe}
\end{equation}
estimated by \eqref{eq:corrected} with the same offset (\cref{prop:pps} holds for any $V$) and a small ridge penalty
$\frac{\lambda}{2}\|\phi\|^2$, $\lambda = 10^{-4}$. Non-hubs have $\phi = 0$, so $\phi_h$ is hub $h$'s
time-invariant attractiveness relative to the rest of the market. It absorbs every fixed attribute of the hub:
unmeasured quality, but also brand, tutorials and default status. Entering the fixed effects as an additional player
in \cref{def:game} gives a share $s_{\mathrm{FE}}$, and
\[
  \bar s_{\mathrm{qual}} = s_{\mathrm{qual}} + s_{\mathrm{FE}}
\]
is the largest share that any time-invariant hub attribute could contribute to quality (\feShare\% of the gain
in-sample).

To see how much of the fixed effects is quality, we project the estimated intercepts on external quality measures
$z_h$ (Open LLM Leaderboard v2 average; arena ratings) with controls $w_h$ (parameter count, re-distributor accounts),
$\hat\phi_h = b_0 + b^\top z_h + \gamma^\top w_h + u_h$, over the \extHubs{} hubs that are covered and precisely
estimated.

\begin{lemma}[Noise-corrected quality share]
\label{lem:noise}
Suppose $\hat\phi_h = \phi_h + e_h$ with $\E[e_h \mid \phi_h, z_h, w_h] = 0$ and $\mathrm{Var}(e_h) = \sigma_h^2$.
Then $\mathrm{Var}(\hat\phi) = \mathrm{Var}(\phi) + \bar\sigma^2$, the projection coefficients are unaffected by $e$,
and
\[
  q = \frac{\mathrm{Var}(\hat b^\top z_h)}{\mathrm{Var}(\hat\phi_h) - \bar\sigma^2}
\]
estimates the share of $\mathrm{Var}(\phi)$ explained by external quality, with $\sigma_h^2$ the bootstrap variance
of $\hat\phi_h$.
\end{lemma}

\noindent We split the fixed-effect share proportionally, $s^{\mathrm{ext}}_{\mathrm{qual}} = s_{\mathrm{qual}} + q\,
s_{\mathrm{FE}}$ ($q = \extQ\%$). Shapley shares are not linear in a split of a player, so this split is an
approximation; the two end points, $s_{\mathrm{qual}}$ and $\bar s_{\mathrm{qual}}$, do not depend on it
(\cref{app:proofs}).

% ---------------------------------------------------------------------------------------------------------------------
\subsection{Does the decomposition recover known mechanisms?}
\label{sec:validation}

The decomposition attributes most of the likelihood gain to popularity (\cref{sec:results}). Because the groups are
correlated, this could in principle be a property of the method rather than of the data. We test the method on
choices whose generating mechanism is known by construction.

For a generating utility $V^{(r)}$, keep every real candidate set $\D_k$ with its real features and inclusion
probabilities, and replace the observed parent $j_k$ by a draw
\begin{equation}
  \tilde j^{(r)}_k \sim P^{(r)}(l \mid \D_k) = \frac{\exp(V^{(r)}_{kl} - \log \pi_{kl})}{\sum_{m \in \D_k}
  \exp(V^{(r)}_{km} - \log \pi_{km})}, \qquad l \in \D_k,
  \label{eq:sir}
\end{equation}
the model-implied choice on the sampled set (\cref{prop:pps}). We use three utilities built from the estimates of
\cref{sec:estimation,sec:bounds}:
\[
  V^{\mathrm{SF}}_{kl} = \hat\theta_{\mathrm{pop}}^\top x_{kl,\mathrm{pop}}, \qquad
  V^{\mathrm{QF}}_{kl} = \hat\theta_{\mathrm{qual}}^\top x_{kl,\mathrm{qual}} + \hat\theta_{\mathrm{match}}^\top
  x_{kl,\mathrm{match}} + \hat q_l, \qquad
  V^{\mathrm{FULL}}_{kl} = \hat V^{\mathrm{FE}}_{kl},
\]
where $\hat q_l = \hat b^\top z_l$ is the part of the hub intercept explained by external quality ($\hat q_l = 0$ for
non-hubs). SF lets choices depend only on popularity, QF only on quality and demand match, and FULL on the estimated
model. Running steps 2--3 of Algorithm 1 on $\{(\tilde j^{(r)}_k, \D_k)\}_k$ gives shares $\hat s(r)$. If the
decomposition separates the mechanisms, $\hat s_{\mathrm{pop}}(\mathrm{SF})$ and $\hat s_{\mathrm{fit}}(\mathrm{QF})$
should be close to one, and $\hat s(\mathrm{FULL})$ should equal the shares estimated on the real choices; $1 -
\hat s_{\mathrm{pop}}(\mathrm{SF})$ and $1 - \hat s_{\mathrm{fit}}(\mathrm{QF})$ measure how much fit leaks to a
correlated group.

The decomposition attributes \recSfPop\% of the gain to popularity under SF and \recQfFit\% to fit under QF, and it
returns \recFullPop\%, \recFullMatch\% and \recFullInertia\% (popularity, demand match, self-supply plus habit) under
FULL, the shares estimated on the real choices.\todo{Recovery is run with four groups; rerun with five or keep the
note that self-supply + habit = the former inertia group.} When choices are generated by quality and demand match
alone, the decomposition does not manufacture a popularity share. As a check of the whole pipeline, we also grow
synthetic networks in which a new node attaches to $l$ with weight $d_l^{\alpha}$ if $d_l > 0$ and $e^{\beta_0}$
otherwise ($\alpha = 1$, $\beta_0 = -1$), replay them with the same code, and recover $\hat\alpha$ within three
standard errors of one under all three sampling designs (\cref{app:validation}).

% ---------------------------------------------------------------------------------------------------------------------
\subsection{Samples, weights and author types}
\label{sec:samples}

The estimator above is defined for any set of occasions and weights. The main specification uses first adoptions
with $w_k = 1$; \cref{sec:results} reports the alternatives defined here, and a rule for calling the results robust
that we fixed before estimating them.

\paragraph{Samples.} \emph{P}: first adoptions (main). \emph{D}: all 24-hour decisions, including repeat decisions.
\emph{U}: all uploads. In D and U the author's earlier parents enter $\D_k$ with $\pi = 1$, so \cref{prop:pps}
applies unchanged. On a repeat decision the degree of a candidate includes the author's own earlier children of it,
which would load habit onto popularity. We therefore replace $d_{kl}$ by $\tilde d_{kl} = d_{kl} - o_{kl}$, where
$o_{kl}$ counts the children of $l$ by author $a_k$ created before $t_k$ (and similarly for adoptions in the last 30
days). On P, $o_{kl} \equiv 0$.

\paragraph{Weights.} With $N$ occasions by $A$ authors and $n_a$ occasions of author $a$, \emph{author weights}
$w_k = (N/A)/n_{a_k}$ give every author the same total weight and a mean weight of one. They estimate the behavior of
the typical author rather than of the typical choice.

\paragraph{Author types.} Authors are classified into six types (major organizations, re-distribution accounts,
other organizations, and individuals with $\ge$100, 10--99 and $<$10 models; \cref{app:data}). Heterogeneity is
tested by $LR = 2\big(\sum_\tau \ell_\tau(\hat\theta_\tau) - \ell(\hat\theta)\big)$, asymptotically $\chi^2$ with
degrees of freedom equal to the number of additional parameters of the type-specific models; because the logit scale may differ across types, types are compared by
shares, not coefficients.

\paragraph{Robustness rule.} An alternative sample or weighting is called robust if its popularity share is within 10
percentage points of the main specification and its fit/popularity ratio within a factor of three.
